%===========================
% File: main.tex
% Arbitrage Theory Course - MMMEF
%===========================
\documentclass[11pt,a4paper,oneside]{book}

%---------------------------
% Encoding & Language
%---------------------------
\usepackage[utf8]{inputenc}
\usepackage[T1]{fontenc}
\usepackage[english]{babel}
\usepackage{lmodern}

%---------------------------
% Mathematics
%---------------------------
\usepackage{amsmath,amssymb,amsthm,mathtools,bbm}
\usepackage{amsfonts}

\usepackage{float}
\numberwithin{equation}{chapter}

%---------------------------
% Figures & Illustrations
%---------------------------
\usepackage{graphicx}
\usepackage{tikz}
\usetikzlibrary{trees, arrows.meta, positioning, calc, shapes.geometric}
\usepackage{tikz-3dplot}

%---------------------------
% Colored Boxes (tcolorbox)
%---------------------------
\usepackage[most]{tcolorbox}

%---------------------------
% Page Layout & Style
%---------------------------
\usepackage{geometry}
\geometry{margin=2.5cm, headheight=14pt, footskip=50pt}

\usepackage{fancyhdr}
\pagestyle{fancy}
\fancyhf{}
\fancyhead[L]{\nouppercase{\leftmark}}
\fancyhead[R]{\thepage}
\fancyfoot[L]{\footnotesize 2025-2026 \textperiodcentered{} Arbitrage Theory - MMMEF}
\fancyfoot[R]{\footnotesize Written by Zakari Hachemi}

\usepackage{enumitem}
\setlist[itemize]{label=\textbullet,leftmargin=1.5em}

\usepackage{array}

%---------------------------
% Hyperlinks
%---------------------------
\usepackage{xcolor}
\usepackage{hyperref}
\usepackage{bookmark}
\usepackage[nameinlink,capitalise,noabbrev]{cleveref}

\hypersetup{
    colorlinks=true,
    linkcolor=blue!50!black,
    citecolor=blue!50!black,
    urlcolor=blue!50!black,
    pdfauthor={Zakari Hachemi},
    pdftitle={MMMEF Course - Arbitrage Theory},
    pdfsubject={Complete Course},
    pdfcreator={LaTeX},
    plainpages=false
}

%---------------------------
% Colors for Environments
%---------------------------
\definecolor{darkblue}{RGB}{0,0,139}
\definecolor{darkgreen}{RGB}{0,100,0}
\definecolor{darkgray}{gray}{0.35}
\definecolor{darkred}{RGB}{99,0,0}

%---------------------------
% Theorem Environments & Colors
%---------------------------

% 1. Standard environment definitions
\theoremstyle{plain}
\newtheorem{theorem}{Theorem}[chapter]
\newtheorem{proposition}[theorem]{Proposition}
\newtheorem{lemma}[theorem]{Lemma}
\newtheorem{corollary}[theorem]{Corollary}

\theoremstyle{definition}
\newtheorem{definition}[theorem]{Definition}
\newtheorem{example}[theorem]{Example}
\newtheorem{assumption}[theorem]{Assumption}

\theoremstyle{remark}
\newtheorem*{remark}{Remark}
\newtheorem*{remarks}{Remarks}

% 2. Apply colors (tcolorbox)

% --- Style for THEOREMS (Red) ---
\tcolorboxenvironment{theorem}{
    enhanced,
    colback=red!5!white,
    colframe=red!75!black,
    boxrule=0.5mm,
    arc=2mm,
    left=2mm, right=2mm, top=2mm, bottom=2mm,
    drop fuzzy shadow
}

% --- Style for PROPOSITIONS (Orange) ---
\tcolorboxenvironment{proposition}{
    enhanced,
    colback=orange!5!white,
    colframe=orange!80!black,
    boxrule=0.5mm,
    arc=2mm,
    left=2mm, right=2mm, top=1mm, bottom=1mm,
    drop fuzzy shadow
}

% --- Style for DEFINITIONS (Blue) ---
\tcolorboxenvironment{definition}{
    enhanced,
    colback=blue!5!white,
    colframe=blue!75!black,
    boxrule=0.5mm,
    arc=2mm,
    left=2mm, right=2mm, top=1mm, bottom=1mm,
    drop fuzzy shadow
}

% --- Style for ASSUMPTIONS (Gray/Blue) ---
\tcolorboxenvironment{assumption}{
    enhanced,
    colback=gray!5!white,
    colframe=gray!75!black,
    boxrule=0.5mm,
    arc=2mm,
    left=2mm, right=2mm, top=1mm, bottom=1mm,
    drop fuzzy shadow
}

% --- Style for EXAMPLES (Green) ---
\tcolorboxenvironment{example}{
    blanker,
    breakable,
    left=5mm,
    before skip=10pt,
    after skip=10pt,
    borderline west={3pt}{0pt}{green!60!black}
}

% --- Style for LEMMAS (Violet) ---
\tcolorboxenvironment{lemma}{
    enhanced,
    colback=violet!5!white,
    colframe=violet!75!black,
    boxrule=0.5mm,
    arc=2mm,
    drop fuzzy shadow
}

% --- Style for COROLLARIES (Cyan) ---
\tcolorboxenvironment{corollary}{
    enhanced,
    colback=cyan!5!white,
    colframe=cyan!75!black,
    boxrule=0.5mm,
    arc=2mm,
    drop fuzzy shadow
}

%---------------------------
% Smart References (cleveref)
%---------------------------
\crefname{theorem}{theorem}{theorems}
\Crefname{theorem}{Theorem}{Theorems}
\crefname{lemma}{lemma}{lemmas}
\Crefname{lemma}{Lemma}{Lemmas}
\crefname{proposition}{proposition}{propositions}
\Crefname{proposition}{Proposition}{Propositions}
\crefname{corollary}{corollary}{corollaries}
\Crefname{corollary}{Corollary}{Corollaries}
\crefname{definition}{definition}{definitions}
\Crefname{definition}{Definition}{Definitions}
\crefname{example}{example}{examples}
\Crefname{example}{Example}{Examples}
\crefname{equation}{equation}{equations}
\Crefname{equation}{Equation}{Equations}
\crefname{chapter}{chapter}{chapters}
\Crefname{chapter}{Chapter}{Chapters}
\crefname{section}{section}{sections}
\Crefname{section}{Section}{Sections}
\crefname{remark}{remark}{remarks}
\Crefname{remark}{Remark}{Remarks}

%---------------------------
% Useful Macros
%---------------------------
\newcommand{\E}{\mathbb{E}}
\newcommand{\PP}{\mathbb{P}}
\newcommand{\RR}{\mathbb{R}}
\newcommand{\indic}{\mathbf{1}}
\DeclarePairedDelimiter{\sca}{\langle}{\rangle}
\newcommand{\circled}[1]{\text{\tikz[baseline=(char.base)]\node[shape=circle,draw,inner sep=0.6pt] (char) {\scriptsize #1};}}

%---------------------------
% Document Information
%---------------------------
\title{\textbf{Arbitrage Theory}\\[0.5em]
       \large Complete Course in English}
\author{Zakari HACHEMI}
\date{Based on the course by Christophe Chorro - MMMEF \\ 2025-2026}

%===========================
\begin{document}
%===========================

\maketitle
\tableofcontents
\mainmatter

%---------------------------
% Course Content
%---------------------------
% ========================================
% Chapter 1: Problem Statement
% ========================================
\chapter{Problem Statement}

\section*{Main References}
\[
\begin{cases}
\text{Mathematics of Arbitrage (Schachermayer \& Delbaen), Chapter 2},\\[3pt]
\text{Introduction to Stochastic Calculus (Lamberton \& Lapeyre), Chapters 1 and 2.}
\end{cases}
\]

\section{The Model}

\begin{definition}[Finite Discrete Time]
Time is modeled by a finite discrete set:
\begin{equation}
t \in \{0, 1, \ldots, T\}.
\end{equation}
\end{definition}

\begin{definition}[Probability Space]
We consider a probability space $(\Omega, \mathcal{F}^{P}, P)$ with $\Omega$ finite, $\mathcal{F}^{P} = \mathcal{P}(\Omega)$ and $\forall \omega \in \Omega,\; P(\omega)>0$. We have $\text{card}(\mathcal{P}(\Omega)) = 2^n$ and $\text{card}(\Omega) = n$.
\end{definition}

\medskip
\begin{tcolorbox}[title=Intuitive Interpretation, colback=yellow!5!white, colframe=yellow!50!black]
\begin{itemize}
    \item $\Omega$: space of possible economic configurations over $[0,T]$;
    \item $\mathcal{F}^{P}$: information available over the period;
    \item $P$: historical probability.
\end{itemize}
\end{tcolorbox}

\begin{definition}[Filtration]
Information is modeled by a filtration $\mathcal{J}^{P} = (\mathcal{J}^{P}_{t})_{t=0,\ldots,T}$ defined by:
\begin{equation}
\mathcal{J}^{P} = (\mathcal{J}^{P}_{t})_{t=0,\ldots,T}, \quad \mathcal{J}^{P}_{0}=\{\emptyset,\Omega\}, \quad \mathcal{J}^{P}_{T}=\mathcal{F}^{P}.
\end{equation}
where $\mathcal{J}^{P}_{0}$ is the smallest $\sigma$-algebra.
\end{definition}

\begin{definition}[Financial Assets]
The market consists of $d+1$ financial assets:
\begin{enumerate}
    \item \textbf{Risk-free asset} (numéraire): $(S_t^0)_{t=0,\ldots,T}$ adapted to $(\mathcal{J}_t)$, with $S_t^0(\omega)>0$.
    \begin{equation}
    (S_t^0)_{t=0,\ldots,T} \text{ adapted to } (\mathcal{J}_t), \quad S_t^0(\omega)>0.
    \end{equation}
    
    \item \textbf{Risky assets}: $d \in \mathbb{N}^*$ assets. The value at date $t$ of asset $j$ is denoted $S_t^j$.
    \begin{equation}
    S_t = (S_t^0, S_t^1, \ldots, S_t^d).
    \end{equation}
\end{enumerate}
\end{definition}

\begin{remarks}
\item The finite $\Omega$ assumption is essential: the spaces $L^0,L^p,L^\infty$ are isomorphic to $\RR^{|\Omega|}$.
\item The risk-free asset plays the role of numéraire: $S_t^0 = e^{rt}$.
\item Measurable random variable:

$X : (\Omega, \mathcal{A}) \longrightarrow (\mathbb{R}, \mathcal{B}(\mathbb{R}))$

$X \text{ measurable random variable} \Longleftrightarrow \forall B \in \mathcal{B}(\mathbb{R}), \; X^{-1}(B) \in \mathcal{A}$

$P_X : \mathcal{B}(\mathbb{R}) \longrightarrow [0,1], \quad B \longmapsto P(X^{-1}(B))$

\vspace{1em}
\textbf{If } $\mathcal{A} = \{\emptyset, \Omega\}$, 
then $X$ random variable $\Longleftrightarrow X$ constant.

\vspace{1em}

\textbf{If } $\mathcal{A} = \mathcal{F}_t = \mathcal{P}(\Omega)$

\text{What is the meaning of being } $\mathcal{F}_t$-measurable? A random variable in this case is simply a mapping.

\vspace{1em}
\item Measurability:
$\hat{S}_{t+1}$ is $\mathcal{F}_{t+1}$-measurable but not $\mathcal{F}_t$-measurable. $\Rightarrow$ One can be in $\mathcal{F}_{t+1}$ without being in $\mathcal{F}_t$

\end{remarks}

\section{Adapted Process to a Filtration}

\begin{tcolorbox}[title=General Intuition, colback=green!5!white, colframe=green!50!black]
A stochastic process $(X_t)_{t=0,\dots,T}$ is said to be \textbf{adapted} to a filtration $(\mathcal{F}_t)_{t=0,\dots,T}$ if, at each time $t$, 
the random variable $X_t$ is \textbf{known at that moment}, meaning it is \textbf{measurable} with respect to $\mathcal{F}_t$.

\medskip
In other words:
\begin{center}
\textbf{Being adapted} $\Longleftrightarrow$ $X_t$ depends only on the information available up to time $t$.
\end{center}
\end{tcolorbox}

\begin{definition}[Formal Definition]
A process $(X_t)_{t=0,\dots,T}$ is \textbf{adapted} to the filtration $(\mathcal{F}_t)$ if:
\[
\forall t \in \{0,1,\dots,T\}, \quad X_t \text{ is } \mathcal{F}_t\text{-measurable.}
\]
In other words:
\[
\forall B \in \mathcal{B}(\mathbb{R}), \quad X_t^{-1}(B) \in \mathcal{F}_t.
\]
\end{definition}

\paragraph*{Interpretation and Importance of Adaptation}
The condition that a price process $(S_t)$ is "adapted" to a filtration $(\mathcal{F}_t)$ is a mathematical formalization of a fundamental and intuitive idea: one cannot know the future. A simple analogy is that of a card game: at each turn $t$, the filtration $\mathcal{F}_t$ represents the set of cards you have seen so far. An "adapted" strategy means that your betting decision at turn $t$ can only depend on these visible cards, and not on cards still in the deck. In finance, this means that the value of an asset $S_t$ at time $t$ is information known at that moment, determined by the history of market events up to $t$, but its future value $S_{t+1}$ remains uncertain.


\begin{example}[Economic Interpretation]
In finance, $\mathcal{F}_t$ represents \textbf{the information available on the market} at time $t$.
Saying that the price of an asset $S_t$ is $\mathcal{F}_t$-measurable means that its value is \textbf{known at date $t$}
from market information.

Thus, an \textbf{adapted} price process $(S_t)$ means that:
\[
S_t \text{ depends on past and present, but not on the future (we don't know } S_{t+1} \text{ at date } t).
\]
\end{example}



\section{Risky Assets and Risk-Free Asset; Reminder on $L^p$ (Finite $\Omega$ Case)}

\begin{definition}[Risky Assets]
Let $d\in\mathbb{N}^*$ be the number of risky assets. For all $j\in\{1,\dots,d\}$ and
$t\in\{0,\dots,T\}$, we denote
\begin{equation}
\hat S_t^{\,j}\quad\text{the value at date }t\text{ of the $j$-th risky asset.}
\end{equation}
We assume that, for all $j$, the process $\big(\hat S_t^{\,j}\big)_{t=0,\dots,T}$ is adapted to the
filtration $\big(\mathcal{J}_t\big)_{t=0,\dots,T}$. \emph{Remark.} Unlike the risk-free asset,
a risky asset can (in principle) reach $0$.
\end{definition}

\begin{definition}[Vector Notation]
We group the prices into a vector
\begin{equation}
\hat S_t:=\big(\hat S_t^{\,0},\hat S_t^{\,1},\dots,\hat S_t^{\,d}\big)\in\RR_+^{\,d+1},
\end{equation}
so that $\hat S=(\hat S_t)_{t=0,\dots,T}$ is a stochastic process with values in $\RR^{d+1}$, adapted to
$(\mathcal{J}_t)$.
\end{definition}

\begin{definition}[Risk-Free Asset and Numéraire]
The risk-free asset (money market account) plays the role of numéraire. In elementary models
(and in this part of the course), we assume that it is possible to borrow/lend at the
constant rate $r>0$ \emph{with continuous compounding}. We then set
\begin{equation}
\hat S_t^{\,0}=e^{rt},\qquad t\in\{0,\dots,T\},
\end{equation}
that is, the value at date $t$ of one euro invested at date $0$.
\end{definition}

\bigskip
\paragraph{Technical Remarks (Finite $\Omega$ Case).}
For reasons of functional analysis, we will assume $\Omega$ \emph{finite}. In particular,
the spaces
\[
L^0(\Omega,\mathcal{F},P),\quad L^p(\Omega,\mathcal{F},P)\ (1\le p<\infty),\quad
L^\infty(\Omega,\mathcal{F},P)
\]
are all canonically identified with $\RR^{|\Omega|}$ (hence a single ``natural'' topology).
We recall the usual definitions:
\[
\begin{aligned}
L^0(\Omega,\mathcal{F},P)
&=\bigl\{\,X:\Omega\to\RR\ \bigm|\ X\ \text{is }\mathcal{F}\text{-measurable}\,\bigr\},\\[2mm]
L^p(\Omega,\mathcal{F},P)
&=\Bigl\{\,X\in L^0\ \Bigm|\ \E\bigl[|X|^p\bigr]<\infty\Bigr\},
\quad 1\le p<\infty,\\[2mm]
L^\infty(\Omega,\mathcal{F},P)
&=\Bigl\{\,X\in L^0\ \Bigm|\ \exists M>0\ \text{such that }P(|X|\le M)=1\Bigr\}.
\end{aligned}
\]
When $\Omega$ is finite, all usual norms are equivalent and
\(
L^0\simeq L^p\simeq L^\infty\simeq \RR^{|\Omega|}.
\)


\subsubsection*{Convention on Interest Rates}

The modeling of the risk-free asset is a cornerstone of financial theory. It serves as a numéraire, that is, a standard of value allowing comparison of monetary flows at different dates. The choice of its dynamics is therefore crucial. In practice, the future value of capital depends on the compounding frequency of interest. The following table illustrates the different conventions and justifies the choice of continuous compounding, $S^0_t = e^{rt}$, which offers mathematical simplification and an idealized representation of a frictionless financial market.

\begin{center}
\begin{tabular}{l|lll}
\textbf{Compounding} & \textbf{Value at $t=1$} & \textbf{Value at $t=T$} \\
\hline
Annual & $(1 + r)$ & $(1 + r)^T$ \\
Semi-annual & $(1 + r/2)^2$ & $(1 + r/2)^{2T}$ \\
Daily & $(1 + r/365)^{365}$ & $(1 + r/365)^{365T}$ \\
$m$ periods/year & $(1 + r/m)^m$ & $(1 + r/m)^{mT}$ \\
\textbf{Continuous ($m \to +\infty$)} & \boldmath{$e^r$} & \boldmath{$e^{rT}$} \\
\end{tabular}
\end{center}

This convention (continuous compounding) is the one we will adopt hereafter: $S^0_t = e^{rt}$.

\section{Binomial Model ($T=3$, $d=1$)}

\subsubsection*{NON RISKY}
$S_0^0 = e^{rt}$

\subsubsection*{RISKY}
$P(w_1, w_2) = P(w_2 \mid w_1) \cdot P(w_1)$

$P(w_1) = \begin{cases}
P(u) \\
P(d)
\end{cases}$

\vspace{1cm}

\begin{center}
\begin{tikzpicture}[
    grow=right,
    scale=0.85, transform shape,
    level distance=45mm,
    level 1/.style={sibling distance=35mm},
    level 2/.style={sibling distance=18mm},
    level 3/.style={sibling distance=9mm},
    edge from parent/.style={draw, -latex},
    every node/.style={align=center}
]

  \node {$s$}
    child {
        node {$sd$}
        child {
            node {$sd^2$}
            child {
                node {$sd^3$}
                edge from parent node[below, font=\footnotesize] {$1-p$}
            }
            child {
                node {$sd^2u$}
                edge from parent node[above, font=\footnotesize] {$p$}
            }
            edge from parent node[below, font=\footnotesize] {$1-p$}
        }
        child {
            node {$sdu$}
            child {
                node {$sd^2u$}
                edge from parent node[below, font=\footnotesize] {$1-p$}
            }
            child {
                node {$sdu^2$}
                edge from parent node[above, font=\footnotesize] {$p$}
            }
            edge from parent node[above, font=\footnotesize] {$p$}
        }
        edge from parent node[below, font=\footnotesize] {$1-p$}
    }
    child {
        node {$su$}
        child {
            node {$sud$}
            child {
                node {$sud^2$}
                edge from parent node[below, font=\footnotesize] {$1-p$}
            }
            child {
                node {$su^2d$}
                edge from parent node[above, font=\footnotesize] {$p$}
            }
            edge from parent node[below, font=\footnotesize] {$1-p$}
        }
        child {
            node {$su^2$}
            child {
                node {$su^2d$}
                edge from parent node[below, font=\footnotesize] {$1-p$}
            }
            child {
                node {$su^3$}
                edge from parent node[above, font=\footnotesize] {$p$}
            }
            edge from parent node[above, font=\footnotesize] {$p$}
        }
        edge from parent node[above, font=\footnotesize] {$p$}
    };
\end{tikzpicture}
\end{center}



\vspace{1cm}

\noindent where

$A_1 = \{\omega \in \Omega \mid \omega_1 = w_u\}$

$A_2 = \{\omega \in \Omega \mid \omega_1 = w_d\}$

$B_1 = \{\omega \in \Omega \mid \omega_1 = w_u, \omega_2 = w_u\}$

$B_2 = \{\omega \in \Omega \mid \omega_1 = w_u, \omega_2 = w_d\}$

$B_3 = \{\omega \in \Omega \mid \omega_1 = w_d, \omega_2 = w_u\}$

$B_4 = \{\omega \in \Omega \mid \omega_1 = w_d, \omega_2 = w_d\}$

\vspace{0.5cm}

$\mathcal{J}_0 = \{\emptyset, \Omega\}, \quad \mathcal{J}_1 = \sigma(A_1, A_2), \quad \mathcal{J}_2 = \sigma(B_1, B_2, B_3, B_4)$

\vspace{1cm}

\noindent \textbf{T=3 Payoff}

\noindent Final Payoff

\section{Investment Strategies}

\begin{assumption}[Frictionless Market Assumption]
We assume the following idealized conditions:
\begin{enumerate}
\item No transaction costs,
\item Possibility of short selling,
\item Possibility of borrowing/lending at the same risk-free rate,
\item Perfect divisibility of assets,
\item No dividends (for simplicity).
\end{enumerate}
\end{assumption}




\begin{tcolorbox}[title=Reminder: Short Selling, colback=blue!5!white, colframe=blue!40!black]
A \textbf{short sale} consists of selling an asset that one does not own, hoping to buy it back later at a lower price.

\medskip
\textbf{Mechanism:}
\begin{enumerate}
    \item The investor ($C_1$) borrows the asset from a third party ($C_2$) through a broker.
    \item He immediately sells this asset on the market at the current price $S_t$.
    \item Later, he buys back the asset at price $S_{t+k}$ to return it (this is the unwinding).
\end{enumerate}

If $S_{t+k} < S_t$, he makes a profit. This is a bearish speculation strategy.

\begin{center}
\begin{tikzpicture}[
  >=Stealth, thick,
  scale=0.8, transform shape,
  circ/.style = {circle, draw, minimum size=14mm, align=center},
  elli/.style = {ellipse, draw, minimum width=30mm, minimum height=16mm, align=center}
]
% Nodes
\node[circ] (c1) {C$_1$};
\node[elli, right=2.8cm of c1] (brk) {Broker};
\node[circ, right=2.8cm of brk] (c2) {C$_2$};

% Arrows
\draw[->] (c1) to[bend left=20] node[midway, above=3pt, font=\small] {short sell} (brk);
\draw[->] (brk) to[bend left=20] node[midway, below=3pt, font=\small] {cash} (c1);

\draw[->] (brk) to[bend left=20] node[midway, above=3pt, font=\small] {borrow security} (c2);
\draw[->] (c2) to[bend left=20] node[midway, below=3pt, font=\small] {security} (brk);
\end{tikzpicture}
\end{center}
\end{tcolorbox}
Breakdown of the operation:
\begin{itemize}
  \item $C_1$ profits if the asset price decreases;
  \item $C_2$ and the broker receive fees;
  \item $C_2$ continues to receive dividends (if any) and can terminate the loan at any time.
\end{itemize}


\section{Financial Strategies and Self-Financing Portfolios}

\begin{definition}[Financial Strategy]
A \textbf{financial strategy} (or investment strategy) is a predictable stochastic process $(\hat H_t)_{t=1,\ldots,T}$ with values in $\RR^{d+1}$:
\begin{equation}
\hat H_t = (\hat H_t^0, \hat H_t^1, \ldots, \hat H_t^d),
\end{equation}
where $\hat H_t^j$ represents the quantity of asset $j$ held in the portfolio over the interval $]t-1, t]$.

The \textbf{portfolio value} (or wealth) at time $t$ is given by:
\begin{equation}
\hat V_t = \langle \hat H_t, \hat S_t \rangle = \sum_{k=0}^d \hat H_t^k \hat S_t^k.
\end{equation}
\end{definition}

\begin{remark}
The previous definition follows from the \emph{frictionless} assumption.
\end{remark}

\begin{definition}[Consumption Process]
Let $(\hat H_t)_{t\in\{0,\ldots,T\}}$ be a financial strategy.  
We define the \textbf{consumption process} $(\hat C_t)_{t\in\{1,\ldots,T\}}$ by:
\[
-\hat C_t = \hat V_t - \langle \hat H_{t-1}, \hat S_t \rangle,
\]
where $\langle \hat H_{t-1}, \hat S_t \rangle$ represents the value of the ``old'' portfolio at time $t$.

A financial strategy is said to be \textbf{self-financing} if, for all $t\in\{1,\ldots,T\}$,
\[
\hat C_t = 0.
\]
\end{definition}

\medskip

In this case, between $0$ and $T$, one does not have the right to consume or invest.  
One can only readjust the portfolio while respecting the budget constraint at each instant.

\begin{example}
If one invests at time $t=0$ and does nothing until maturity $T$, one implements a \emph{self-financing} strategy.
\end{example}

\begin{lemma}
A financial strategy can be characterized by $(\hat V_0, (\hat C_t)_{t\in\{1,\ldots,T\}}, (\hat H_t)_{t\in\{0,\ldots,T\}})$.
\end{lemma}

\begin{remark}
The lemma formalizes the link between financial flows (buying, selling, consumption) and portfolio composition.
The idea is that at each period $t$, we want to know:
\begin{itemize}
    \item how much of each asset we have (the vector $\hat H_t$);
    \item how much the portfolio is worth ($\hat V_t$);
    \item how much we withdraw or add (consumption $\hat C_t$ or capital injection).
\end{itemize}
The lemma states that the data of the initial value, the consumption process (how much is taken out of the portfolio at each date) and the risky strategy (how risk is allocated) is sufficient to completely reconstruct the strategy, including the portion invested in the risk-free asset (which serves as an adjustment or financing account).
\end{remark}


\begin{proof}
We denote $\hat S_t=(\hat S_t^0,\hat S_t^1,\ldots,\hat S_t^d)$ the price vector at time $t$ (non-discounted),
$\hat H_t=(\hat H_t^0,\ldots,\hat H_t^d)$ the strategy, and
\[
\hat V_t \;=\; \langle \hat H_t,\hat S_t\rangle \;=\; \sum_{k=0}^d \hat H_t^k\,\hat S_t^k
\]
the portfolio value at time $t$.
The \emph{consumption process} $(\hat C_t)_{t=1,\ldots,T}$ is defined by
\begin{equation}\label{eq:consumption-def}
-\hat C_t \;=\; \hat V_t - \langle \hat H_{t-1}, \hat S_t\rangle
\qquad (t=1,\ldots,T),
\end{equation}
that is: we start from the value of the old portfolio after the price movement,
we consume $\hat C_t$ (positive = withdrawal, negative = contribution), then we reallocate to obtain $\hat H_t$.

\smallskip
\noindent\textbf{Step 1 ($t=0$).}
By definition,
\[
\hat V_0 \;=\; \sum_{k=0}^d \hat H_0^k\,\hat S_0^k
\;=\; \hat H_0^0\,\hat S_0^0 \;+\; \sum_{k=1}^d \hat H_0^k\,\hat S_0^k.
\]
Isolating the risk-free component, we immediately obtain
\[
\hat H_0^0 \;=\; \frac{\hat V_0 - \sum_{k=1}^d \hat S_0^k\,\hat H_0^k}{\hat S_0^0}.
\]

\smallskip
\noindent\textbf{Step 2 ($t\ge 1$).}
The budget constraint at time $t$ is obtained by rewriting \eqref{eq:consumption-def}:
\[
\hat V_t \;=\; \langle \hat H_{t-1}, \hat S_t\rangle - \hat C_t.
\]
Now $\hat V_t=\langle \hat H_t,\hat S_t\rangle = \hat H_t^0\,\hat S_t^0 + \sum_{k=1}^d \hat H_t^k\,\hat S_t^k$,
so that
\[
\hat H_t^0\,\hat S_t^0 \;+\; \sum_{k=1}^d \hat H_t^k\,\hat S_t^k
\;=\; \langle \hat H_{t-1}, \hat S_t\rangle \;-\; \hat C_t.
\]
Isolating $\hat H_t^0$, we obtain the announced formula:
\[
\hat H_t^0 \;=\; \frac{-\hat C_t + \langle \hat H_{t-1}, \hat S_t\rangle - \sum_{k=1}^d \hat H_t^k\,\hat S_t^k}{\hat S_t^0},
\qquad t=1,\ldots,T.
\]

\smallskip
\noindent\textbf{Conclusion (characterization).}
The identities above show that, once we fix
$\hat V_0$, the consumption process $(\hat C_t)$ and the risky components
$(\hat H_t^1,\ldots,\hat H_t^d)$ for each $t$, the risk-free component $\hat H_t^0$
is entirely determined for $t=0$ and then recursively for $t\ge 1$.
Thus, the triplet $(\hat V_0,(\hat C_t),(\hat H_t))$ completely characterizes the financial strategy.
\end{proof}

\begin{corollary}
A self-financing financial strategy can be characterized by $(\hat V_0, (\hat H_t)_{t\in\{0,\ldots,T\}})$.
\end{corollary}

\begin{definition}[Discounted Values]
For all $k\in\{0,\ldots,d\}$, we define the \textbf{discounted values} by:
\[
S_t^k = \frac{\hat S_t^k}{\hat S_t^0}, \qquad V_t = \frac{\hat V_t}{\hat S_t^0}.
\]
Thus, $S_t^0 \equiv 1$ and $S_t = (S_t^1,\ldots,S_t^d)$ represents the vector of risky assets only.
\end{definition}

\begin{lemma}

\emph{A financial strategy is self-financing}
\[
\Longleftrightarrow\;
\forall\, t \in \{1,\ldots,T\},\quad
V_t = V_0 + \sum_{k=0}^{t-1} \langle H_k,\, S_{k+1} - S_k \rangle.
\]
\end{lemma}



\begin{proof}
By definition, for all $t\in\{0,\ldots,T\}$:
\[
\hat V_t = \hat H_t^0 \hat S_t^0 + \sum_{k=1}^d \hat H_t^k \hat S_t^k = \langle \hat H_t, \hat S_t \rangle.
\]

The self-financing condition gives:
\[
\hat V_t = \hat H_{t-1}^0 \hat S_t^0 + \sum_{k=1}^d \hat H_{t-1}^k \hat S_t^k = \hat H_{t-1}^0 \hat S_t^0 + \langle \hat H_{t-1}, \hat S_t \rangle.
\]
Thus,
\[
V_t = V_{t-1} + \langle H_{t-1}, S_t - S_{t-1} \rangle,
\]
and by induction:
\[
V_t = V_0 + \sum_{k=0}^{t-1} \langle H_k, S_{k+1} - S_k \rangle.
\]
\end{proof}

\begin{remark}
The variation in portfolio value is entirely attributable to the variation in asset prices when the strategy is self-financing (i.e., without consumption).
\end{remark}

\begin{remark}
We denote:
\[
(H \cdot S)_t = \sum_{k=0}^{t-1} \langle H_k, S_{k+1} - S_k \rangle,
\]
which corresponds to the discrete stochastic integral of process $H$ with respect to $S$.
\end{remark}


\section{Contingent Claims and Attainability Sets}

A contingent claim (or derivative asset) is a financial contract whose value at a future date $T$ depends on the future state of the world. Mathematically, it is a random variable $\hat{U}$ that is $\mathcal{F}_T$-measurable. Its discounted value is denoted $U = \hat{U} / S^0_T$. The central question of the theory is whether one can "attain" this claim with a given initial cost.

\begin{itemize}
    \item \textbf{Replicable claim:} A claim $\hat{U}$ is said to be replicable at price $a$ if there exists a self-financing financial strategy $(H_t)$ such that the initial investment is $\hat{V}_0 = a$ and the final portfolio value exactly matches the claim's payoff, $\hat{V}_T = \hat{U}$.
    \item \textbf{Set $K_a$:} This set represents all final payoffs (discounted) that can be attained with an initial cost $a$. Formally:
    \[
    K_a = \{ a + (H \cdot S)_T \mid H \in \mathcal{H}^d \}
    \]
    where $\mathcal{H}^d$ is the set of predictable strategies on the risky assets. The set $K = K_0$ is particularly important: it represents the set of net gains attainable at zero initial cost.
    \item \textbf{Set $\mathcal{C}_a$:} This set represents all claims (discounted) that can be super-replicated with an initial cost $a$, meaning their payoff is dominated by an attainable payoff:
    \[
    \mathcal{C}_a = \{ U \in L^\infty(\mathcal{F}_T) \mid \exists g \in K_a, g \ge U \}
    \]
    The set $\mathcal{C} = \mathcal{C}_0$ corresponds to claims super-replicable at zero cost.
    \item \textbf{Properties of $K$ and $\mathcal{C}$:} These sets have fundamental geometric structures for what follows.
    \begin{itemize}
        \item $K$ is a vector space (closed in finite dimension).
        \item $\mathcal{C}$ is a convex cone (closed in finite dimension).
    \end{itemize}
    These geometric properties are fundamental because they validate the application of Hahn-Banach separation theorems, which constitute the cornerstone of the proof of the Fundamental Theorem of Arbitrage.
\end{itemize}

\section{Arbitrage Opportunities}

\begin{definition}[Arbitrage Opportunity]
A self-financing strategy $(\hat H_t)_{t\in[0,T]}$ is an \emph{arbitrage opportunity} (A.O. hereafter) if
\begin{equation}
\hat V_0 = 0, \quad \hat V_T \ge 0 \ \text{P-a.s.}, \quad \text{and} \quad P(\hat V_T > 0) > 0.
\end{equation}
\end{definition}

\begin{remarks}
\item[(a)] In the previous definition, the condition above remains unchanged if the probability $P$ is replaced by a probability $Q$ such that $Q \sim P$.
\item[(b)] The following conditions are equivalent:
\begin{align}
& V_0 = 0, \; V_T \ge 0 \text{ P-a.s. and } P(V_T > 0) > 0 \\
\Longleftrightarrow \quad & V_0 = 0, \; (H \cdot S)_T \ge 0 \text{ P-a.s. and } P((H \cdot S)_T > 0) > 0 \\
\Longleftrightarrow \quad & V_0 = 0, \; V_T \ge 0 \text{ P-a.s. and } \E_P[V_T] > 0.
\end{align}
\end{remarks}

\paragraph*{Context of Attainability Sets}
The central objective of arbitrage theory is to determine whether a future "payoff", that is, a contingent claim such as an option, can be perfectly generated by a dynamic investment strategy. To formalize this idea, we introduce two fundamental sets. The set $K_a$ represents the universe of all possible financial outcomes (expressed in discounted value at the final date $T$) that can be attained starting from an initial capital $a$ and applying a self-financing strategy. The set $\mathcal{C}_a$, on the other hand, is broader: it groups all contingent claims that can at least be matched, that is, "super-replicated", with this same initial capital $a$. These geometric structures are the key to linking the absence of risk-free profit opportunities to the valuation of financial assets.

\begin{definition}[Contingent Claim]
\begin{enumerate}[label=(\alph*)]
\item A contingent claim $\hat U$ is a random variable measurable with respect to $\mathcal{F}_T$.  
In this case, we define its discounted value by
\begin{equation}
U = \frac{\hat U}{\hat S_T^0}.
\end{equation}

\item A contingent claim $\hat U$ is said to be \emph{replicable} (resp. \emph{super-replicable}) at price $a$ if there exists a self-financing portfolio $(\hat H_t)_{t\in[0,T]}$ such that
\begin{equation}
\hat V_0 = a \quad \text{and} \quad \hat V_T = \hat U \; (\text{resp. } \hat V_T \ge \hat U).
\end{equation}
We then denote
\begin{equation}
K_a = \bigl\{\, a + (H \cdot S)_T \;|\; H \in \mathcal{H}^d \,\bigr\},
\end{equation}
where $\mathcal{H}^d$ is the set of adapted stochastic processes with values in $\mathbb{R}^d$, and
\begin{equation}
\mathcal{C}_a = \bigl\{\, U \in L^\infty(\mathcal{F}_T) \;|\; \exists g \in K_a,\, g \ge U \,\bigr\}.
\end{equation}
\end{enumerate}
We always have $K_a \subset \mathcal{C}_a$.
\end{definition}

\begin{remarks}
\item[(a)] We denote $\mathcal{C} = \mathcal{C}_0$ and $K = K_0$. Thus, $\mathcal{C}_a = a + \mathcal{C}$ and $K_a = a + K$.

\item[(b)] $K$ is a vector space (hence convex and closed in finite dimension).

\item[(c)] $\mathcal{C}$ is a convex cone containing
\[
L_-^\infty(\Omega, \mathcal{F}_T, P) = \{\, X \in L^\infty \mid X \le 0 \text{ P-a.s.} \}.
\]
\end{remarks}

\begin{definition}[No-Arbitrage Assumption (NA)]
We say that the \emph{no-arbitrage} assumption (NA hereafter) is satisfied if there exists no arbitrage opportunity, that is:
\begin{equation}
\text{if } V_0 = 0, \; V_T \ge 0 \text{ P-a.s. } \Rightarrow P(V_T = 0) = 1.
\end{equation}
\end{definition}

\begin{tcolorbox}[title=Intuition: No Free Lunch, colback=red!5!white, colframe=red!50!black]
The NA assumption means that it is impossible to make money for sure without initial investment and without taking any risk of loss.
\begin{itemize}
    \item If I start from nothing ($V_0=0$) and can never lose ($V_T \ge 0$), then I cannot gain anything ($V_T=0$).
    \item This is the minimal condition for a financial model to be viable. Without it, prices would have no meaning (infinite supply/demand).
\end{itemize}
\end{tcolorbox}

\begin{lemma}
We have the following equivalences:
\begin{equation}
NA \;\Longleftrightarrow\; K \cap L_+^0(\Omega, \mathcal{F}_T, P) = \{0\}
\;\Longleftrightarrow\;
\mathcal{C} \cap L_+^0(\Omega, \mathcal{F}_T, P) = \{0\}.
\end{equation}
\end{lemma}

\begin{proof}
Direction ($NA \Rightarrow K \cap L^0_+ = \{0\}$): We reason by contradiction. Suppose that the no-arbitrage assumption (NA) is satisfied, but that there exists a non-zero element $g$ in the intersection $K \cap L^0_+$.
\begin{enumerate}
    \item Since $g \in K$, by definition, $g$ is the discounted final value $V_T$ of a self-financing strategy $H$ with initial cost $V_0 = 0$. So $g = (H \cdot S)_T$.
    \item Since $g \in L^0_+$, we have $g \ge 0$ P-a.s. (almost surely).
    \item Since $g$ is non-zero, $P(g > 0) > 0$.
    \item In summary, we have found a self-financing strategy with $V_0 = 0$, $V_T \ge 0$ P-a.s., and $P(V_T > 0) > 0$. This is precisely the definition of an arbitrage opportunity. However, we assumed NA. There is therefore a contradiction. The initial hypothesis is false, and therefore $K \cap L^0_+$ must be reduced to $\{0\}$.
\end{enumerate}

Direction ($K \cap L^0_+ = \{0\} \Rightarrow NA$): Suppose that $K \cap L^0_+ = \{0\}$. We must show that the NA assumption is satisfied.
\begin{enumerate}
    \item Consider any self-financing strategy $H$ such that $V_0 = 0$ and $V_T \ge 0$ P-a.s.
    \item By definition, the final value $V_T$ of such a strategy belongs to $K$.
    \item Since $V_T \ge 0$, $V_T$ also belongs to $L^0_+$.
    \item Therefore, $V_T$ belongs to the intersection $K \cap L^0_+$.
    \item By our hypothesis, this intersection is reduced to $\{0\}$. This implies that $V_T = 0$ P-a.s. This corresponds exactly to the definition of the no-arbitrage assumption (NA): any strategy starting from zero and that can only generate gains cannot actually generate any gain.
\end{enumerate}

This mathematical reformulation of the no-arbitrage condition is the key that will allow us to mobilize the tools of functional analysis to prove the Fundamental Theorem of asset pricing.
\end{proof}

\begin{definition}[Local Arbitrage]
Let $t \in \{1, \dots, T\}$.  
A \emph{local arbitrage at $t$} is a random variable $\xi$ measurable with respect to $\mathcal{F}_{t-1}$, with values in $\mathbb{R}^d$, such that:
\[
\langle \xi, S_t - S_{t-1} \rangle \ge 0 \text{ P-a.s.}, 
\quad \text{and} \quad P\bigl( \langle \xi, S_t - S_{t-1} \rangle > 0 \bigr) > 0.
\]
The market satisfies condition $NA(t)$ if there exists no local arbitrage at $t$.
\end{definition}

\begin{lemma}
\[
NA \;\Longleftrightarrow\; NA(t) \quad \forall t \in \{1, \dots, T\}.
\]
\end{lemma}

\begin{proof}[Proof of the lemma: $NA \Longleftrightarrow NA(t)\ \forall t$]
\subsection*{Direction $\Rightarrow$:}
Suppose that the global no-arbitrage assumption (NA) is satisfied.  
If there existed a local arbitrage at some $t \in \{1,\ldots,T\}$, then one could construct a self-financing strategy $(H_s)_{s\in[0,T]}$ such that
\[
V_0 = 0, \qquad 
V_T = \langle \xi, S_t - S_{t-1} \rangle \ge 0 \text{ a.s.}, \qquad 
\mathbb{P}(V_T > 0) > 0.
\]
This strategy would constitute a global arbitrage opportunity, which would contradict (NA).  
Thus, if (NA) is true, no local arbitrage can exist:
\[
NA \;\Rightarrow\; NA(t), \quad \forall t\in\{1,\ldots,T\}.
\]

\subsection*{Direction $\Leftarrow$:}
We show the converse by induction on the number of periods $T$ of the market.

\subsubsection*{Base Case}
For $T=1$, the result is obvious: a global arbitrage over a single period is exactly a local arbitrage.

\subsubsection*{Induction Hypothesis}
Suppose the result is true for a market of maturity $T-1$, that is:
\[
\text{in a market with $T-1$ periods,} \qquad
NA \;\Leftarrow\; NA(t), \quad \forall t \in \{1,\ldots,T-1\}.
\]

\subsubsection*{Induction Step}
Consider a market with $T$ periods and reason by contraposition.  
Let $(\hat H_t)_{t=0}^T$ be an arbitrage opportunity (A.O.) on $[0,T]$.  
We then have:
\[
V_T = V_{T-1} + \langle H_{T-1}, S_T - S_{T-1} \rangle, 
\qquad V_T \ge 0,\quad \mathbb{P}(V_T > 0) > 0.
\]
Let us introduce the set and the random variable:
\[
A = \{ V_{T-1} < 0 \} \in \mathcal{F}_{T-1},
\qquad
\xi = H_{T-1}\mathbf{1}_A \quad (\mathcal{F}_{T-1}\text{-measurable}).
\]
Then:
\[
\mathbf{1}_A V_T
= \mathbf{1}_A V_{T-1}
+ \big\langle \mathbf{1}_A H_{T-1},\, S_T - S_{T-1} \big\rangle
\ge 0.
\]

\textbf{Case 1:} If $\mathbb{P}(A) > 0$, then $\xi$ defines a local arbitrage because
\[
\langle \xi, S_T - S_{T-1} \rangle \ge 0 \text{ a.s.}, 
\qquad 
\mathbb{P}(\langle \xi, S_T - S_{T-1} \rangle > 0) \ge \mathbb{P}(A) > 0.
\]

\textbf{Case 2:} If $\mathbb{P}(A) = 0$, then $V_{T-1} \ge 0$ a.s.  
Two sub-cases arise:
\[
\begin{cases}
(1) & V_{T-1} = 0\ \text{a.s.} 
\quad \Rightarrow\ H_{T-1}\ \text{is a local arbitrage at } T, \\[0.5em]
(2) & \mathbb{P}(V_{T-1}>0) > 0 
\quad \Rightarrow\ (H_t)_{t\le T-1}\ \text{is an A.O. on } [0,T-1].
\end{cases}
\]
In the second sub-case, the induction hypothesis applies: there exists a local arbitrage 
at some $t \in \{1,\ldots,T-1\}$.

Thus, in all cases, a global A.O. on $[0,T]$ implies the existence of a local arbitrage.  
We deduce:
\[
NA(t)\ \forall t \;\Rightarrow\; NA.
\]

\textbf{Conclusion:} both directions are thus established, and we indeed have
\[
NA \;\Longleftrightarrow\; NA(t), \quad \forall t\in\{1,\ldots,T\}.
\]
\end{proof}

% ========================================
% Chapter 2: Fundamental Theorem of Asset Pricing (FTAP)
% ========================================
\chapter{Fundamental Theorem of Asset Pricing (FTAP)}

Fundamental Theorem of Asset Pricing

\section{Reference}
\textbf{Harrison--Pliska (1981)}

\begin{theorem}[Fundamental Theorem of Asset Pricing (FTAP)]
The no-arbitrage assumption is equivalent to the existence of an equivalent martingale measure:
\begin{equation}
NA \Longleftrightarrow \mathcal{M}^e(S) \neq \varnothing,
\end{equation}
where $\mathcal{M}^e(S)$ denotes the set of \textbf{equivalent martingale measures}:
\begin{equation}
\mathcal{M}^e(S) = \left\{\, Q \sim P \ \Big|\ (S_t)_{t=0,\dots,T} \text{ is a } Q\text{-martingale} \right\}.
\end{equation}
In other words, $Q \in \mathcal{M}^e(S)$ if and only if:
\begin{enumerate}
    \item $Q$ is a probability measure equivalent to $P$ (i.e., $Q(\omega) > 0$ for all $\omega \in \Omega$);
    \item The discounted price process $(S_t)$ (with $S^0 \equiv 1$) is a martingale under $Q$.
\end{enumerate}
\end{theorem}

\section{Definition and Properties of the Set $\Delta$}

\begin{definition}
We define:
\begin{equation}
\Delta = \Bigl\{\, Y \in L_+^0(\Omega, \mathcal{J}, P) \ \big|\ \mathbb{E}_P[Y] = 1 \Bigr\}.
\end{equation}
\end{definition}

\begin{lemma}
The set $\Delta$ is convex and compact.
\end{lemma}

\begin{proof}
\begin{enumerate}
    \item \textbf{Convexity:} Let $Y_1, Y_2 \in \Delta$ and $\lambda \in [0,1]$. We have $Y_1 \ge 0$ and $Y_2 \ge 0$, so $\lambda Y_1 + (1-\lambda)Y_2 \ge 0$. Moreover, by linearity of expectation, $\E_P[\lambda Y_1 + (1-\lambda)Y_2] = \lambda \E_P[Y_1] + (1-\lambda)\E_P[Y_2] = \lambda(1) + (1-\lambda)(1) = 1$. Therefore $\lambda Y_1 + (1-\lambda)Y_2 \in \Delta$. The set is indeed convex.
    \item \textbf{Closedness:} The set $\Delta$ can be seen as the intersection of two sets:
    \begin{itemize}
        \item The cone of positive random variables $L^0_+$, which is a closed set.
        \item The affine hyperplane $\{Y \in L^0 \mid \E_P[Y] = 1\}$. This set is the preimage of the closed singleton $\{1\}$ under the linear (and hence continuous in finite dimension) mapping $\varphi(Y) = \E_P[Y]$. It is therefore closed.
    \end{itemize}
    The intersection of two closed sets is a closed set.
    \item \textbf{Boundedness:} For all $Y \in \Delta$, we have by definition $Y \ge 0$ and $\E_P[Y] = 1$. The expectation reads $\sum_{\omega \in \Omega} Y(\omega)P(\omega) = 1$. Since each term in the sum is positive, we have for any $\omega$:
    \[
    Y(\omega)P(\omega) \le \sum_{\omega' \in \Omega} Y(\omega')P(\omega') = 1
    \]
    Since $P(\omega) > 0$ for all $\omega$ (model assumption), we can divide to obtain:
    \[
    0 \le Y(\omega) \le \frac{1}{P(\omega)}
    \]
    Since the universe $\Omega$ is finite, there exists a finite upper bound for all $Y(\omega)$: $\max_{\omega \in \Omega} (1 / P(\omega))$. The set $\Delta$ is therefore bounded.
    \item \textbf{Conclusion:} In the space $\mathbb{R}^{|\Omega|}$ (to which $L^0$ is isomorphic), the set $\Delta$ is closed and bounded. By the Heine-Borel theorem, it is therefore compact.
\end{enumerate}
\end{proof}

\begin{tcolorbox}[title=Proof Strategy for FTAP (Direction $\Rightarrow$), colback=orange!5!white, colframe=orange!50!black]
The proof of the implication $NA \Rightarrow \mathcal{M}^e(S) \neq \varnothing$ is the heart of the Fundamental Theorem. It is structured in five successive logical steps:

\begin{enumerate}
    \item \textbf{Geometric reformulation (Lemma 2.4):} We translate the economic condition of no arbitrage into a geometric condition: the sets $K$ (attainable gains at zero cost) and $\Delta$ (probability densities) must be disjoint.
    \item \textbf{Separation (Lemma 2.6):} We apply the Hahn-Banach separation theorem (geometric form) to these two disjoint convex sets, one closed ($K$) and the other compact ($\Delta$). This guarantees the existence of a separating hyperplane.
    \item \textbf{Construction of the density $Z$ (Lemma 2.7):} We exploit the properties of the separating hyperplane to construct a strictly positive random variable $Z$ such that $\E_P[Zk] = 0$ for all $k \in K$.
    \item \textbf{Construction of measure $Q$ (Corollary 2.8):} We use this density $Z$ to define a new probability measure $Q$, equivalent to $P$, under which all gains from $K$ have zero expectation.
    \item \textbf{Martingale Property (Lemma 2.9):} We finally show that if all attainable gains have zero expectation under $Q$, then the price process $S$ is a $Q$-martingale.
\end{enumerate}
\end{tcolorbox}

The following sections detail each of these steps.

\section{Link Between Absence of Arbitrage and the Set $\Delta$}

\begin{lemma}
\begin{equation}
NA \Longleftrightarrow K \cap \Delta = \varnothing.
\end{equation}
\end{lemma}

\begin{proof}
We know that:
\[
NA \Longleftrightarrow K \cap L_+^0 = \{0\}.
\]
We must therefore show that $K \cap L_+^0 = \{0\}$ is equivalent to $K \cap \Delta = \varnothing$.

\medskip
\noindent
($\Rightarrow$)  
If $K \cap L_+^0 = \{0\}$, then necessarily $K \cap \Delta = \varnothing$.

\medskip
\noindent
($\Leftarrow$)  
By contraposition, suppose there exists $g \in K \cap L_+^0$ with $g \neq 0$.  
We then define:
\[
Y = \frac{g}{\mathbb{E}_P[g]}.
\]
We indeed have $\mathbb{E}_P[Y] = 1$ and $Y \ge 0$, so $Y \in \Delta$.  
Since $Y \in K$ as well (because $K$ is a vector space and $Y$ is an element of $K$ multiplied by a scalar), this contradicts $K \cap \Delta = \varnothing$.  
Thus, the two conditions are equivalent.
\end{proof}

\section{Vector Spaces and Continuous Linear Functionals}

\paragraph{Notations.}
If $E$ and $F$ are vector spaces, we denote:
\[
\mathcal{L}_c(E,F)
\]
the set of continuous linear maps from $E$ to $F$.

\medskip
\noindent
Recall that for any vector space $E$, 
\[
\dim(E) = \dim(\ker(\Phi)) + \operatorname{rank}(\Phi).
\]
Thus, if $\Phi \in \mathcal{L}_c(E,F)$ is non-zero, then $\operatorname{rank}(\Phi)=1$, and consequently:
\[
\dim(\ker(\Phi)) = \dim(E) - 1.
\]
The kernel of $\Phi$ is therefore a \textbf{hyperplane} of $E$.

\medskip
\noindent
In our setting, if $E = L^0(\Omega, \mathcal{J}, P)$ and $F = \mathbb{R}$, then $E \simeq \mathbb{R}^{|\Omega|}$ (since $\Omega$ is finite).  
Any element $\Phi \in \mathcal{L}_c(E,F)$ can be expressed using the values taken on indicator functions.

For any $X \in L^0(\Omega, \mathcal{J}, P)$, we have:
\[
X(\omega) = \sum_{\omega' \in \Omega} X(\omega') \mathbf{1}_{\{\omega'\}}(\omega),
\qquad \text{hence} \qquad
X = \sum_{\omega' \in \Omega} X(\omega') \mathbf{1}_{\{\omega'\}}.
\]
Thus, for any continuous linear functional $\Phi$, we can write:
\[
\Phi(X) = \sum_{\omega' \in \Omega} X(\omega') \, \Phi(\mathbf{1}_{\{\omega'\}}).
\]

In particular, if we write:
\[
Z(\omega) = \frac{\Phi(\mathbf{1}_{\{\omega\}})}{P(\omega)},
\]
then
\[
\Phi(X) = \sum_{\omega' \in \Omega} X(\omega') \frac{\Phi(\mathbf{1}_{\{\omega'\}})}{P(\omega')} P(\omega')
= \mathbb{E}_P[X Z].
\]
Thus, any continuous linear functional on $L^0$ can be represented as a \textbf{weighted expectation}.


\section{Separation Lemma}

\begin{theorem}[Separation Theorem]
In a Banach space, if $A$ and $B$ are two convex, closed, and disjoint sets, one of which is compact,  
then there exists $\Phi \in \mathcal{L}_c(\text{Banach}, \mathbb{R})$ and $\alpha \in \mathbb{R}$ such that:
\[
\Phi(a) < \alpha,\ \forall a \in A, \qquad \Phi(b) > \alpha,\ \forall b \in B.
\]
In our case, we apply this theorem with $A = K$ and $B = \Delta$.
\end{theorem}

\begin{lemma}[Separation Lemma] \label{lem:6}
If $NA$ is satisfied, then there exists $\Phi \in \mathcal{L}_c(L^0(\Omega, \mathcal{J}, P), \mathbb{R})$ and $\alpha \in \mathbb{R}$ such that:
\begin{equation}
\forall k \in K,\ \Phi(k) < \alpha, \qquad \forall y \in \Delta,\ \Phi(y) > \alpha.
\end{equation}
\end{lemma}

\begin{remark}[Key Remarks]
\item This is called the separation theorem because we can separate A and B with an affine hyperplane.
\item Since $K$ is a vector subspace, if $\mathbb{E}_P[Zk] \neq 0$ for some $k \in K$,  
then for $t \to \pm\infty$, we would have $tk \in K$ and:
\[
\mathbb{E}_P[Z(tk)] = t \, \mathbb{E}_P[Zk],
\]
which would violate the strict bound $\Phi(k) < \alpha$.  
Thus, necessarily:
\[
\mathbb{E}_P[Zk] = 0,\quad \forall k \in K.
\]
\end{remark}


\begin{lemma}\label{lem:7}
Under the No-Arbitrage assumption (NA), there exists
\[
 Z\in L^0(\Omega,\mathcal{J},P)\simeq \RR^{|\Omega|}
\]
such that
\begin{equation}
\forall k\in K,\quad \E_P[Zk]=0
\qquad\text{and}\qquad
\forall \omega\in\Omega,\ Z(\omega)>0 .
\end{equation}
\end{lemma}

\begin{proof}
Using the Lemma and the previous remark, there exists
\[
 Z\in L^0(\Omega,\mathcal{J},P)\ \ (\text{isomorphic to } \RR^{|\Omega|})
 \qquad\text{and}\qquad
 \alpha\in\RR
\]
such that
\[
\forall k\in K,\quad \E_P[Zk]<\alpha
\qquad\text{and}\qquad
\forall Y\in\Delta,\ \E_P[YZ]>\alpha .
\]
Taking $k=0$ we deduce $\alpha>0$.

Then define the linear functional
\[
\varphi:K\longrightarrow\RR,\qquad
k\longmapsto \E_P[Zk] .
\]
This is a linear map defined on the vector space $K$ and it is bounded.
By separation, we must have $\varphi\equiv 0$, i.e.,
\[
\forall k\in K,\qquad \E_P[Zk]=0.
\]

Finally, for $\omega_0\in\Omega$, set
\[
Y=\frac{\mathbbm{1}_{\{\omega_0\}}}{P(\omega_0)}\in\Delta.
\]
Then
\[
\E_P[YZ]>\alpha>0
\qquad\Longrightarrow\qquad
Z(\omega_0)=\E_P[YZ]>0 .
\]
Since $\omega_0$ is arbitrary, we obtain $Z(\omega)>0$ for all $\omega\in\Omega$.
\end{proof}

\begin{corollary}\label{cor:Q}
NA $\Rightarrow$ There exists a probability $Q\sim P$ such that
\[
\forall k\in K,\qquad \E_Q[k]=0 .
\]
\end{corollary}

\begin{proof}[Proof of Corollary \ref{cor:Q}]
Let $Z$ be given by Lemma 7 and set
\[
\widetilde Z=\frac{Z}{\E_P[Z]},
\qquad
Q(\omega)=\widetilde Z(\omega)\,P(\omega).
\]
Since $Z\ge0$ and $\E_P[\widetilde Z]=1$, $Q$ is a probability and $Q\sim P$.
Moreover, for all $k\in K$,
\[
\E_Q[k]=\E_P[\widetilde Z\,k]
=\frac{1}{\E_P[Z]}\,\E_P[Zk]=0 .
\]
\end{proof}

\begin{lemma}\label{lem:8}
NA $\Rightarrow\ \mathcal{M}^e(S)\neq\varnothing$. 
\end{lemma}

\begin{proof}[Detailed proof of Lemma \ref{lem:8}]
Let $Q$ be the probability measure whose existence is guaranteed by Corollary \ref{cor:Q}. This measure satisfies $\E_Q[k]=0$ for all $k \in K$. We will show that under this measure, the current price process $(S_t)$ is a martingale.

\begin{enumerate}
    \item \textbf{Objective:} We must show that for all assets $j \in \{1,\dots,d\}$ and all times $t \in \{0,\dots,T-1\}$, the conditional expectation of price increments is zero:
    \[
    \E_Q[S_{t+1}^j - S_t^j \mid \mathcal{J}_t] = 0.
    \]
    
    \item \textbf{Reformulation:} By definition of conditional expectation, this amounts to showing that for all events $A \in \mathcal{J}_t$ (the information available at $t$), the following expectation is zero:
    \[
    \E_Q[\mathbbm{1}_A (S_{t+1}^j - S_t^j)] = 0.
    \]
    
    \item \textbf{Construction of the Test Strategy:} To prove this point, we construct an ad-hoc financial strategy $H$ that "bets" on this asset over this time interval if event $A$ occurs. Define the predictable process $H = (H_u)$ by:
    \[
    H_u^k = 
    \begin{cases} 
    \mathbbm{1}_A & \text{if } k=j \text{ and } u=t, \\
    0 & \text{otherwise.}
    \end{cases}
    \]
    This strategy consists of holding one unit of asset $j$ between $t$ and $t+1$ if and only if $A$ is realized. It is indeed $\mathcal{J}_u$-measurable for all $u$ (since $A \in \mathcal{J}_t$ so $H_t$ is $\mathcal{J}_t$-measurable).
    
    \item \textbf{Link with Space $K$:} The final gain generated by this strategy (starting from zero capital) is:
    \[
    (H \cdot S)_T = \sum_{u=0}^{T-1} \langle H_u, S_{u+1} - S_u \rangle = \langle H_t, S_{t+1} - S_t \rangle = \mathbbm{1}_A (S_{t+1}^j - S_t^j).
    \]
    By definition, this gain belongs to the space $K$ of attainable gains at zero cost. Therefore, $\mathbbm{1}_A (S_{t+1}^j - S_t^j) \in K$.
    
    \item \textbf{Conclusion:} According to Corollary \ref{cor:Q}, we know that $\E_Q[k] = 0$ for all $k \in K$. Applying this to our specific gain, we obtain:
    \[
    \E_Q[\mathbbm{1}_A (S_{t+1}^j - S_t^j)] = 0.
    \]
    This being true for all $A \in \mathcal{J}_t$, all $t$ and all $j$, we have demonstrated that $(S_t)$ is a $Q$-martingale. Consequently, $Q \in \mathcal{M}^e(S)$, and the set of equivalent martingale measures is non-empty.
\end{enumerate}
\end{proof}

\begin{lemma}\label{lem:9}
If $Q\in\mathcal{M}^e(S)$, then for any strategy $H\in\mathcal{H}^d$,
the process $\big((H\cdot S)_t\big)_{t=0,\dots,T}$ is a martingale under $Q$.
\end{lemma}

\begin{proof}[Proof of Lemma \ref{lem:9}]
We know that
\[
(H\cdot S)_{t+1}=(H\cdot S)_t+\big\langle H_t,\,S_{t+1}-S_t\big\rangle .
\]
Taking the conditional expectation with respect to $\mathcal{J}_t$,
\[
\E_Q\!\big[(H\cdot S)_{t+1}\mid\mathcal{J}_t\big]
=(H\cdot S)_t+\Big\langle H_t,\ \E_Q\!\big[S_{t+1}-S_t\mid\mathcal{J}_t\big]\Big\rangle .
\]
Under $Q$, $(S_t)$ is a martingale, so $\E_Q[S_{t+1}-S_t\mid\mathcal{J}_t]=0$ and thus
\[
\E_Q\!\big[(H\cdot S)_{t+1}\mid\mathcal{J}_t\big]=(H\cdot S)_t .
\]
\end{proof}

\begin{theorem}[Fundamental Theorem of Asset Pricing — "Thm 1"]
\begin{equation}
NA\ \Longleftrightarrow\ \mathcal{M}^e(S)\neq\varnothing .
\end{equation}
\end{theorem}

\begin{tcolorbox}[title=Intuition: The Economic Meaning of FTAP, colback=blue!5!white, colframe=blue!50!black]
The Fundamental Theorem of Asset Pricing establishes a bridge between an economic condition (absence of arbitrage) and a mathematical concept (existence of a risk-neutral probability). But what is the deep intuition?

\begin{itemize}
    \item \textbf{A World Without Expected Return:} Under a risk-neutral measure $Q$, the expected return of all risky assets, once discounted, is exactly equal to the return of the risk-free asset (here 0 since discounted). In other words, their expected "excess return" is zero.
    \item \textbf{No Strategy Can Win on Average:} An investment strategy is a dynamic combination of these assets. If each base component has zero expected return (relative to the risk-free asset), then no combination of these components can have a strictly positive expected return if it starts from zero initial capital.
    \item \textbf{The Impossibility of the "Money Machine":} An arbitrage opportunity is precisely a strategy that starts from zero, cannot lose, and has a chance of winning. The existence of a measure $Q$ eliminates this possibility by guaranteeing that any strategy starting from zero has zero expected gain. It transforms a viable market into a "fair game" from the perspective of expected returns.
\end{itemize}
\end{tcolorbox}

\begin{proof}
\emph{$\Rightarrow$}: given by Lemmas \ref{lem:8} and \ref{lem:7} (via Corollary \ref{cor:Q}).

\smallskip
\noindent
(\emph{$\Leftarrow$}): $\mathcal{M}^e(S) \neq \varnothing \implies NA$.

\smallskip
Suppose that $\mathcal{M}^e(S) \neq \varnothing$ and let $Q \in \mathcal{M}^e(S)$ be an equivalent martingale measure.
We will show that the condition $K \cap L_+^0 = \{0\}$ is satisfied, which is equivalent to $NA$.

Let $g \in K \cap L_+^0$.
\begin{enumerate}
    \item By definition of $K$, there exists a self-financing strategy $H$ such that $V_0 = 0$ and $V_T = (H \cdot S)_T = g$.
    \item Since $Q \in \mathcal{M}^e(S)$, the price process $S$ is a $Q$-martingale. According to Lemma \ref{lem:9}, the value process $(H \cdot S)_t$ is also a $Q$-martingale (since $H$ is predictable and bounded in a finite space).
    \item The martingale property implies that the expectation of the final gain equals the initial value:
    \[
    \E_Q[g] = \E_Q[V_T] = V_0 = 0.
    \]
    \item By hypothesis, $g \in L_+^0$, that is, $g \ge 0$ $P$-a.s.
    Since $Q$ is equivalent to $P$ ($Q \sim P$), the set of null measure events is identical for both measures. Therefore, $g \ge 0$ $Q$-a.s.
    \item We have a random variable $g$ such that $g \ge 0$ $Q$-a.s. and $\E_Q[g] = 0$.
    Now, the integral of a positive function is zero only if the function is zero almost everywhere.
    We thus deduce that $g = 0$ $Q$-a.s.
    \item By equivalence of measures, $g = 0$ $P$-a.s.
\end{enumerate}
Thus, the only element of $K \cap L_+^0$ is the zero element. The intersection is reduced to $\{0\}$, which proves absence of arbitrage ($NA$).
\end{proof}

% ========================================
% Chapter 3: Complete Markets
% ========================================
\chapter{Complete Markets}

\begin{definition}
The market is \emph{complete} if every bounded contingent claim is replicable, that is:
\begin{equation}
\forall \hat U \in L^\infty(\Omega,\mathcal{J}_T,P),\ \exists a\in\RR,\ \exists H\in\mathcal{H}^d
\text{ such that } U = a + (H\cdot S)_T.
\end{equation}
\end{definition}

\begin{remark}
Under (NA), if a claim $\hat U$ is replicable at price $a$, then $a$ and the trajectory $(H\cdot S)_t$ are unique, and:
\[
\forall Q\in\mathcal{M}^e(S),\quad
a = \E_Q[U], \qquad a + (H\cdot S)_t = \E_Q[U \mid \mathcal{J}_t].
\]
\end{remark}

\begin{proof}[Proof of the remark]
\textbf{Uniqueness of price $a$.}  
Suppose that $U=a_1+(H^1\cdot S)_T=a_2+(H^2\cdot S)_T$ with $a_1>a_2$.  
Then
\[
\big((H^1 - H^2)\cdot S\big)_T = a_2 - a_1 < 0,
\]
and the reverse strategy generates a positive certain gain — an arbitrage opportunity, contradiction.  
Therefore $a_1=a_2$.

\smallskip
\textbf{Uniqueness of strategy.}  
Suppose $U=a+(H^1\cdot S)_T=a+(H^2\cdot S)_T$ but $(H^1\cdot S)\neq(H^2\cdot S)$.  
There then exists $t$ and a set $A\in\mathcal{J}_t$ such that
\[
A = \{(H^1\cdot S)_t > (H^2\cdot S)_t\},\quad P(A)>0.
\]
Define the strategy $H = (H^1-H^2)\indic_A \indic_{[t,T]}$.  
Then $(H\cdot S)_T = \indic_A((H^1-H^2)\cdot S)_T - \indic_A((H^1-H^2)\cdot S)_t \ge 0$ with strictly positive probability, hence arbitrage. Contradiction.

\smallskip
\textbf{Expectation relations under $Q$.}  
According to the previous lemma, if $Q\in\mathcal{M}^e(S)$, then $(H\cdot S)_t$ is a martingale under $Q$, hence:
\[
a = \E_Q[U],\qquad a + (H\cdot S)_t = \E_Q[U\mid\mathcal{J}_t].
\]
\end{proof}

\begin{theorem}[Characterization of Complete Markets]
\label{thm:complet_singleton}
Under (NA), the market is complete if and only if the set of equivalent martingale measures is a singleton:
\begin{equation}
\text{Complete} \ \Longleftrightarrow\ \mathcal{M}^e(S)=\{Q\}.
\end{equation}
\end{theorem}

\begin{tcolorbox}[title=Summary: Completeness and Uniqueness, colback=purple!5!white, colframe=purple!50!black]
This result is central:
\begin{itemize}
    \item \textbf{Incomplete Market} ($\mathcal{M}^e(S)$ infinite): There exists intrinsic unhedgeable risk. The price of an option is not unique (price interval).
    \item \textbf{Complete Market} ($\mathcal{M}^e(S) = \{Q\}$): Every risk can be perfectly hedged. The price is unique and given by the expectation under the unique measure $Q$.
\end{itemize}
In the binomial model (Chap. 1 and 5), the market is complete.
\end{tcolorbox}

\begin{lemma}\label{lem:10}
\label{lem:UinK_iff_EQ0}
Let $U\in L^\infty(\Omega,\mathcal{J}_T,P)$. Under (NA),
\begin{equation}
U\in K \quad \Longleftrightarrow \quad \forall Q\in\mathcal{M}^e(S),\ \E_Q[U]=0.
\end{equation}
\end{lemma}

\begin{proof}
This lemma is a powerful tool that characterizes the space $K$ of attainable gains without initial cost.

\emph{Direction $\Rightarrow$.} If $U\in K$, then $U = (H\cdot S)_T$ for some strategy $H$ with zero initial cost $\hat{V}_0=0$. For all $Q\in\mathcal{M}^e(S)$, the process $((H\cdot S)_t)$ is a $Q$-martingale. Consequently, using the martingale property at $t=0$ and $t=T$:
\[
\E_Q[U] = \E_Q[(H\cdot S)_T] = (H\cdot S)_0 = \hat{V}_0 = 0.
\]

\emph{Direction $\Leftarrow$.} Suppose that $\E_Q[U] = 0$ for all $Q \in \mathcal{M}^e(S)$, but that $U \notin K$.
We will use a separation argument, similar to the FTAP proof.
\begin{enumerate}
    \item The sets $\{U\}$ (which is convex and compact since reduced to a point) and $K$ (which is a vector space, hence convex and closed) are disjoint.
    \item By the Hahn-Banach separation theorem, there exists a continuous linear form $\lambda$ (identifiable with a random variable since $\Omega$ is finite) such that:
    \[
    \langle \lambda, U \rangle > 0 \quad \text{and} \quad \langle \lambda, k \rangle = 0 \quad \forall k \in K.
    \]
    The second condition implies $\E_P[\lambda k] = 0$ for all $k \in K$.
    \item The problem is that $\lambda$ is not necessarily positive everywhere, so it does not directly define a probability measure.
    \item Let $\tilde{Q} \in \mathcal{M}^e(S)$ (whose existence is guaranteed under NA). Since $\Omega$ is finite and $\tilde{Q}(\omega)>0$, we can choose $\varepsilon > 0$ small enough so that the random variable $Z' = \frac{d\tilde{Q}}{dP} + \varepsilon \lambda$ is strictly positive everywhere:
    \[
    \tilde{Q}(\omega) + \varepsilon \lambda(\omega) P(\omega) > 0 \quad \forall \omega.
    \]
    \item We can then define a new probability $\bar{Q}$ by its density (up to normalization) proportional to $Z'$.
    \item Let us verify that $\bar{Q}$ is a martingale measure. For all $k \in K$:
    \[
    \E_{\bar{Q}}[k] \propto \E_P[k \cdot (\frac{d\tilde{Q}}{dP} + \varepsilon \lambda)] = \E_{\tilde{Q}}[k] + \varepsilon \E_P[k \lambda] = 0 + 0 = 0.
    \]
    Therefore $\bar{Q} \in \mathcal{M}^e(S)$.
    \item Now compute the expectation of $U$ under $\bar{Q}$:
    \[
    \E_{\bar{Q}}[U] \propto \E_{\tilde{Q}}[U] + \varepsilon \E_P[U \lambda].
    \]
    Now, $\E_{\tilde{Q}}[U]=0$ by hypothesis (since $\tilde{Q} \in \mathcal{M}^e(S)$) and $\E_P[U \lambda] = \langle \lambda, U \rangle > 0$ by separation.
    Therefore $\E_{\bar{Q}}[U] > 0$.
    \item We have constructed a martingale measure $\bar{Q} \in \mathcal{M}^e(S)$ for which $\E_{\bar{Q}}[U] \neq 0$, which contradicts our initial hypothesis. Therefore, $U$ must belong to $K$.
\end{enumerate}
\end{proof}

\subsection{The Case of Non-Replicable Assets}
When the market is incomplete, the set $\mathcal{M}^e(S)$ contains more than one martingale measure. In this case, not all assets are replicable, and their price, if it exists, is no longer unique. The following result precisely formalizes the structure of possible prices.

\begin{proposition}[No-Arbitrage Price Interval]
Under the no-arbitrage assumption (NA) and for a contingent claim $U$:
\begin{enumerate}
    \item[(a)] If $U$ is replicable, then $\bar{\Pi}(U) = \Pi(U)$. The set of no-arbitrage prices for $U$ is the singleton $\{\Pi(U)\}$, where $\Pi(U)$ is the unique value $\E_Q[U]$ valid for all $Q \in \mathcal{M}^e(S)$.
    \item[(b)] If $U$ is not replicable, then $\bar{\Pi}(U) > \underline{\Pi}(U)$. The set of no-arbitrage prices for $U$ is the open interval $]\underline{\Pi}(U), \bar{\Pi}(U)[$.
\end{enumerate}
\end{proposition}

\paragraph*{Intuition: The No-Arbitrage "Bid-Ask Spread"}
Market incompleteness means that there exists risk that cannot be entirely eliminated by hedging. There is no longer a unique consensus on valuation, but a range of "legitimate" prices. This range can be interpreted as a theoretical bid-ask spread.
\begin{itemize}
    \item The upper bound, $\bar{\Pi}(U) = \sup \E_Q[U]$, represents the \textbf{ask price}. It is the minimum cost at which a seller is willing to guarantee delivery of payoff $U$ (via a super-replication strategy), which corresponds to the maximum price the buyer will have to pay.
    \item The lower bound, $\underline{\Pi}(U) = \inf \E_Q[U]$, represents the \textbf{bid price}. It is the maximum amount a buyer is willing to pay to acquire the claim, which corresponds to the minimum price the seller can hope to receive.
\end{itemize}
Each measure $Q$ in $\mathcal{M}^e(S)$ represents a pricing model consistent with absence of arbitrage. The multiplicity of these measures translates directly into a range of possible prices for non-replicable assets.

\begin{proof}[Proof of Theorem~\ref{thm:complet_singleton}]
This theorem links the economic notion of completeness to the mathematical structure of the set of martingale measures.

\emph{Direction $\Rightarrow$ (Complete $\Rightarrow$ Singleton)}:
Suppose the market is complete. Let $Q_1$ and $Q_2$ be any two measures in $\mathcal{M}^e(S)$.
For any event $A \in \mathcal{F}_T$, the contingent claim $U = \mathbf{1}_A$ is bounded.
Since the market is complete, $U$ is replicable. There therefore exists a price $a$ and a strategy $H$ such that $\mathbf{1}_A = a + (H \cdot S)_T$.
Taking expectation under $Q_1$ and $Q_2$ (knowing that the expectation of the stochastic part is zero under a martingale measure):
\begin{itemize}
    \item $\E_{Q_1}[\mathbf{1}_A] = \E_{Q_1}[a + (H \cdot S)_T] = a + 0 = a$
    \item $\E_{Q_2}[\mathbf{1}_A] = \E_{Q_2}[a + (H \cdot S)_T] = a + 0 = a$
\end{itemize}
We thus have $Q_1(A) = a$ and $Q_2(A) = a$.
Since $Q_1(A) = Q_2(A)$ for all events $A$, the measures are identical: $Q_1 = Q_2$. The set $\mathcal{M}^e(S)$ can therefore contain only one element.

\smallskip
\emph{Direction $\Leftarrow$ (Singleton $\Rightarrow$ Complete)}:
Suppose that $\mathcal{M}^e(S) = \{Q\}$. Let $U$ be any bounded contingent claim. We want to show that it is replicable.
\begin{enumerate}
    \item Set $a = \E_Q[U]$. This is the candidate price.
    \item Consider the "centered" claim $U - a$. For any measure $Q'$ in $\mathcal{M}^e(S)$, we must have $\E_{Q'}[U-a] = 0$.
    \item Since there is only one measure $Q$, this condition reduces to $\E_Q[U-a] = \E_Q[U] - a = 0$, which is true by definition of $a$.
    \item According to Lemma \ref{lem:UinK_iff_EQ0}, since $U-a$ has zero expectation under all measures in $\mathcal{M}^e(S)$, we have $U-a \in K$.
    \item By definition of $K$, there exists a strategy $H$ such that $U-a = (H \cdot S)_T$.
    \item This is the definition of replicability: $U = a + (H \cdot S)_T$. The claim $U$ is replicable at price $a$. The market is therefore complete.
\end{enumerate}
\end{proof}


\section{Pricing by Absence of Arbitrage (Kreps, 1981): Static Approach}

Let $U\in L^\infty(\Omega,\mathcal{J}_T,P)$ be a claim (we work with discounted values).  
We extend the market by adding a synthetic asset $U$: at $t=0$, it can be bought/sold at price $a\in\RR$, and it pays $U$ at $T$ (no intermediate payments).

An extended strategy is written $(\theta,\bar H)$ where $\theta\in\RR$ is the position on asset $U$ and $\bar H\in\mathcal{H}^{d}$ is the strategy on the $d$ risky assets.  
If $\bar H$ is self-financing, then $(\theta,\bar H)$ is self-financing and
\begin{equation}
V_T = V_0 + \theta\,(U-a) + (H\cdot S)_T.
\end{equation}
We define
\begin{equation}
K_{a,U} := \{\ \theta\,(U-a) + k \ :\ \theta\in\RR,\ k\in K\ \}
= \RR\,(U-a)+K,
\end{equation}
the set of attainable gains at zero initial cost in the extended market.

\begin{definition}
A real number $a$ is a \emph{no-arbitrage price} for $U$ if the extended market satisfies (NA), i.e.,
\begin{equation}
K_{a,U}\cap L_+^0(\Omega,\mathcal{J}_T,P)=\{0\}.
\end{equation}
\end{definition}

\begin{remark}
From the \emph{FTAP} and the characterization $Q\in\mathcal{M}^e(S)\iff \forall k\in K,\ \E_Q[k]=0$, we deduce:
\[
a \text{ is no-arbitrage } \ \Longleftrightarrow\ \exists Q\in\mathcal{M}^e(S)\ \text{ such that }\ \E_Q[U-a]=0
\ \Longleftrightarrow\ \exists Q\in\mathcal{M}^e(S)\ \text{ with } a=\E_Q[U].
\]
\end{remark}

\begin{theorem}[Envelope of No-Arbitrage Prices]
\label{thm:intervalle_prix_NA}
Under (NA), the set of no-arbitrage prices for $U\in L^\infty$ is the open interval (non-empty and bounded)
\begin{equation}
\big(\, \inf_{Q\in\mathcal{M}^e(S)} \E_Q[U]\ ,\ \sup_{Q\in\mathcal{M}^e(S)} \E_Q[U] \,\big),
\end{equation}
where the bounds are given by the extreme values of $\E_Q[U]$ as $Q$ ranges over $\mathcal{M}^e(S)$.
In particular, if the market is complete, this set reduces to the unique value $a=\E_Q[U]$, where $Q$ is the unique equivalent martingale measure.
\end{theorem}

\begin{proof}[Proof Idea]
($\subset$) If $a$ is no-arbitrage, there exists $Q\in\mathcal{M}^e(S)$ with $a=\E_Q[U]$ (previous remark).  
($\supset$) If $a=\E_Q[U]$ for some $Q\in\mathcal{M}^e(S)$, then for all $k'\in K_{a,U}$,
\[
\E_Q[k'] = \E_Q[\theta(U-a)+k]= \theta\,\E_Q[U-a]+\E_Q[k]=0,
\]
so the extended market satisfies (NA) by FTAP.  
The lower/upper bound results from compactness of $\mathcal{M}^e(S)$ on finite space.
\end{proof}



\begin{remark}\label{rq:UinK_implique_0}
If $U\in K$, then for any equivalent martingale measure $Q\in\mathcal{M}^e(S)$ we have
\[
\E_Q[U]=0.
\]
In particular, the set of possible no-arbitrage prices for $U$ is reduced to $\{0\}$.
\end{remark}

\begin{definition}[No-Arbitrage Bounds]
For $U\in L^\infty(\Omega,\mathcal{J}_T,P)$, we define
\begin{equation}
\bar\Pi(U)=\sup\{\E_Q[U]\ :\ Q\in\mathcal{M}^e(S)\}, \qquad
\underline{\Pi}(U)=\inf\{\E_Q[U]\ :\ Q\in\mathcal{M}^e(S)\}.
\end{equation}
\end{definition}

\begin{lemma}\label{lem:replicabilite_intervalle_prix}
Under (NA) and for a claim $U\in L^\infty$:
\begin{enumerate}
\item[(a)] If $\bar\Pi(U)=\underline{\Pi}(U)$, then $U$ is replicable at the unique price
\[
\Pi(U)=\bar\Pi(U)=\underline{\Pi}(U),
\]
and the set of no-arbitrage prices is the singleton $\{\Pi(U)\}$.
\item[(b)] If $\bar\Pi(U)>\underline{\Pi}(U)$, then $U$ is not replicable and
the set of no-arbitrage prices is the open interval
\[
\big(\,\underline{\Pi}(U),\,\bar\Pi(U)\,\big).
\]
\end{enumerate}
\end{lemma}

\begin{proof}
(a) The set $\{\E_Q[U]:Q\in\mathcal{M}^e(S)\}$ is non-empty (FTAP) and compact (finite space).
If it reduces to a single point $m$, set $f=U-m$. Then $\E_Q[f]=0$ for all $Q\in\mathcal{M}^e(S)$,

Therefore $U=m+f$ is replicable at price $m=\bar\Pi(U)=\underline{\Pi}(U)$ and this price is unique
(\emph{cf.} uniqueness remark above).

(b) If $U$ were replicable, there would exist $a$ such that $\E_Q[U]=a$ for all $Q\in\mathcal{M}^e(S)$,
which would force $\bar\Pi(U)=\underline{\Pi}(U)=a$, contradiction. Therefore $U$ is not replicable.

Moreover, the image of the convex set $\mathcal{M}^e(S)$ by the affine map $Q\mapsto \E_Q[U]$
is an interval of $\RR$; we already know it is contained in $[\underline{\Pi}(U),\bar\Pi(U)]$,
and it is not reduced to a point. It remains to see that the extremities are not attained by a no-arbitrage price
in the \emph{extended} market.

Suppose for example that there exists $Q^\star\in\mathcal{M}^e(S)$ such that
$\E_{Q^\star}[U]=\bar\Pi(U)$. By compactness, we can find $\tilde Q\in\mathcal{M}^e(S)$
with $\E_{\tilde Q}[U]<\E_{Q^\star}[U]$. For small $\varepsilon>0$, define
$Q_\varepsilon \propto Q^\star+\varepsilon(Q^\star-\tilde Q)$; for $\varepsilon$ small enough,
$Q_\varepsilon\in\mathcal{M}^e(S)$ (same argument as in the proof of EMM existence),
and
\[
\E_{Q_\varepsilon}[U]
=\frac{\E_{Q^\star}[U]+\varepsilon\big(\E_{Q^\star}[U]-\E_{\tilde Q}[U]\big)}
      {\sum_\omega \big(Q^\star(\omega)+\varepsilon(Q^\star-\tilde Q)(\omega)\big)}
>\E_{Q^\star}[U],
\]
contradiction with the definition of $\bar\Pi(U)$. An identical argument applies for $\underline{\Pi}(U)$.
We deduce that the set of no-arbitrage prices is exactly the open interval
$\big(\underline{\Pi}(U),\bar\Pi(U)\big)$.
\end{proof}

\section{Dynamic Approach: No-Arbitrage Price Process}

\begin{definition}\label{def:processus_prix_NA}
A process $(A_t)_{t=0..T}$ is a \emph{no-arbitrage price process} for a claim $U$
if $A_T=U$ and if the market extended to $(S^1,\dots,S^d,A)$ (still in discounted values)
satisfies (NA).
\end{definition}

\begin{lemma}\label{lem:prix_dynamique_EQ}
If the initial market $(S^1,\dots,S^d)$ satisfies (NA), then $(A_t)$ is a no-arbitrage price process
for $U$ if and only if there exists $Q\in\mathcal{M}^e(S)$ such that
\[
A_t=\E_Q[U\mid \mathcal{J}_t],\qquad t=0,\dots,T.
\]
\end{lemma}

\begin{proof}
This is a direct application of Definition \ref{def:processus_prix_NA} and the FTAP:
in the extended market, (NA) $\Leftrightarrow$ existence of a $Q$ under which all discounted assets
(including $A$) are martingales, hence $A_t=\E_Q[A_T\mid\mathcal{J}_t]=\E_Q[U\mid\mathcal{J}_t]$.
The converse is immediate.
\end{proof}

\begin{lemma}\label{lem:unicite_processus_si_replicable}
Suppose (NA) and $U$ is replicable. Then:
\begin{enumerate}
\item[(i)] The no-arbitrage price process for $U$ is unique.
\item[(ii)] For all $Q\in\mathcal{M}^e(S)$, we have $A_t=\E_Q[U\mid\mathcal{J}_t]$.
\item[(iii)] A process $(H_t)$ is a replication process for $U$ if and only if
\[
A_0=\bar\Pi(U)\quad\text{and}\quad A_t=\bar\Pi(U)+(H\cdot S)_t,\ \ t=0,\dots,T.
\]
\end{enumerate}
\end{lemma}

\begin{proof}
If $U=\bar\Pi(U)+(H\cdot S)_T$ is replicable, then for all $Q\in\mathcal{M}^e(S)$,
\[
A_t:=\E_Q[U\mid\mathcal{J}_t]=\bar\Pi(U)+\E_Q[(H\cdot S)_T\mid\mathcal{J}_t]
=\bar\Pi(U)+(H\cdot S)_t,
\]
since $(H\cdot S)$ is a martingale under $Q$. This expression depends neither on the choice
of $Q$ nor on another process $A$, hence (i) and (ii). Point (iii) is the direct equivalence:
$A_T=U$ if and only if $A_t=\bar\Pi(U)+(H\cdot S)_t$ for all $t$ (martingale property).
\end{proof}

\section{Super-Replication and Super-Hedging Bound}

\begin{definition}
We say that $U$ is \emph{super-replicable} at price $a\in\RR$ if there exists $H\in\mathcal{H}^d$ such that
\[
a + (H\cdot S)_T \ge U\ \ \text{$P$-a.s.}
\]
We set
\begin{equation}
A(U)=\big\{\, a\in\RR\ :\ \exists H\in\mathcal{H}^d,\ a+(H\cdot S)_T\ge U \,\big\},
\qquad
\alpha := \inf A(U).
\end{equation}
\end{definition}

\begin{theorem}\label{thm:superhedging_equals_upper_envelope}
Under (NA), we have $A(U)=[\alpha,+\infty)$ and
\begin{equation}
\alpha=\bar\Pi(U)=\sup_{Q\in\mathcal{M}^e(S)} \E_Q[U].
\end{equation}
\end{theorem}

\begin{proof}
Theorem 3.12 establishes a remarkable duality result: this cost, obtained by a primal approach (portfolio search), equals the supremum of expected prices under all martingale measures (dual approach).

\textbf{Step 1: Show that $A(U)$ is a closed interval $[\alpha, +\infty)$.}
\begin{itemize}
    \item If $a \in A(U)$, there exists $k \in K$ such that $a+k \ge U$. If $b > a$, then $b+k = (a+k) + (b-a) \ge U + (b-a) > U$. We can find a portfolio that generates $b+k$, so $b \in A(U)$. The set is an interval.
    \item To show it is closed, we must prove that the set $K + L^0_+$ is closed. In finite dimension ($\Omega$ finite), this is the case. Let $f_m = k_m + l_m$ be a sequence converging to $f$, with $k_m \in K$ and $l_m \in L^0_+$. Let $Q \in \mathcal{M}^e(S)$. By continuity of expectation, $\E_Q[f_m] \to \E_Q[f]$. The sequence $(\E_Q[f_m])$ is therefore bounded. Now, $\E_Q[f_m] = \E_Q[k_m] + \E_Q[l_m] = \E_Q[l_m]$. Thus, the sequence $(\E_Q[l_m])$ is bounded. Since $l_m(\omega) \ge 0$ and $Q(\omega) > 0$ for all $\omega$ ($\Omega$ finite), this implies that the sequence $(l_m)$ is bounded in $\mathbb{R}^{|\Omega|}$. By the Bolzano-Weierstrass theorem, we can extract a convergent subsequence $l_{\phi(m)} \to l \in L^0_+$. Consequently, $k_{\phi(m)} = f_{\phi(m)} - l_{\phi(m)}$ converges to $f-l$. Since $K$ is closed, $f-l \in K$. Finally, $f = (f-l) + l \in K + L^0_+$, proving the set is closed.
\end{itemize}

\textbf{Step 2: Show that $\alpha \ge \bar\Pi(U)$.}
\begin{itemize}
    \item Let $a \in A(U)$ be any super-replication cost. By definition, there exists $k \in K$ (representing the gain of a zero-cost portfolio) such that $a + k \ge U$.
    \item Take the expectation of this inequality under any measure $Q \in \mathcal{M}^e(S)$.
    \item $\E_Q[a + k] \ge \E_Q[U]$. By linearity, $a + \E_Q[k] \ge \E_Q[U]$.
    \item Since $k \in K$, we know that $\E_Q[k]=0$. The inequality becomes $a \ge \E_Q[U]$.
    \item This being true for all measures $Q \in \mathcal{M}^e(S)$, we deduce that $a$ is greater than or equal to the supremum: $a \ge \sup_{Q \in \mathcal{M}^e(S)} \E_Q[U] = \bar\Pi(U)$.
    \item Finally, since this holds for all costs $a \in A(U)$, it also holds for their infimum $\alpha$. Therefore, $\alpha \ge \bar\Pi(U)$.
\end{itemize}

\textbf{Step 3: Show that $\alpha \le \bar\Pi(U)$.}
\begin{itemize}
    \item We reason by contradiction assuming $\alpha > \bar\Pi(U)$. We can then choose a real $a$ such that $\bar\Pi(U) < a < \alpha$. By definition of $\alpha$, $a$ is not a super-replication cost, so $a \notin A(U)$.
    \item This means that the set $\{a - U + k \mid k \in K\}$ is disjoint from the set of positive variables $L^0_+$.
    \item We have two disjoint closed convex sets. The separation theorem gives us the existence of a linear form $\lambda \ge 0$ ($\lambda \neq 0$) such that $\langle k, \lambda \rangle = 0$ for all $k \in K$ and $\langle a - U, \lambda \rangle \le 0$.
    \item Normalizing $\lambda$, we construct a probability $Q = \lambda / \sum \lambda(\omega)$. This probability $Q$ is a martingale measure ($\E_Q[k]=0$), but not necessarily equivalent to $P$. The separation inequality implies $a \le \E_Q[U]$.
    \item If $Q$ were equivalent ($Q \in \mathcal{M}^e(S)$), we would have $a \le \E_Q[U] \le \bar\Pi(U)$, contradicting $a > \bar\Pi(U)$.
    \item If $Q$ is not equivalent, we must refine the argument. Take a measure $\tilde{Q} \in \mathcal{M}^e(S)$ (which exists under NA). For all $\varepsilon \in (0, 1)$, define the measure $Q_\varepsilon = \varepsilon \tilde{Q} + (1-\varepsilon)Q$. This is a convex combination of martingale measures, so $Q_\varepsilon$ is also a martingale measure. Moreover, $Q_\varepsilon$ is equivalent to $P$ because $\tilde{Q}$ is. Therefore $Q_\varepsilon \in \mathcal{M}^e(S)$.
    \item By definition of $\bar\Pi(U)$, we have $\bar\Pi(U) \ge \E_{Q_\varepsilon}[U]$ for all $\varepsilon$.
    \item Developing: $\bar\Pi(U) \ge \varepsilon \E_{\tilde{Q}}[U] + (1-\varepsilon) \E_Q[U]$.
    \item Since $a \le \E_Q[U]$, we have $\bar\Pi(U) \ge \varepsilon \E_{\tilde{Q}}[U] + (1-\varepsilon)a$.
    \item This inequality holds for all $\varepsilon \in (0, 1)$. Taking $\varepsilon$ to 0, we obtain $\bar\Pi(U) \ge a$.
    \item This contradicts our initial hypothesis $a > \bar\Pi(U)$. The hypothesis $\alpha > \bar\Pi(U)$ is therefore false. We must have $\alpha \le \bar\Pi(U)$.
\end{itemize}

The two inequalities combined prove that $\alpha = \bar\Pi(U)$.
\end{proof}

% ========================================
% Chapter 4: Examples of One-Period Models
% ========================================
\chapter{Examples of One-Period Models}

\section{The Binomial Model}
We consider a one-period market $T=1$ with a universe $\Omega = \{\omega_u, \omega_d\}$. The historical probability is $P(\omega_u)=p > 0$ and $P(\omega_d)=1-p$.
\begin{itemize}
    \item A risk-free asset whose discounted price is constant: $S_0^0 = 1, S_1^0 = e^r$. The interest rate is $r > 0$.
    \item A risky asset with initial price $S_0^1=s$ and final prices $S_1^1(\omega_u) = su$ and $S_1^1(\omega_d) = sd$, with $u > d$.
\end{itemize}

\begin{proposition}
The no-arbitrage assumption (NA) is satisfied if and only if $d < e^r < u$.
In this case, there exists a unique equivalent martingale probability measure $Q \in \mathcal{M}^e(S)$, defined by:
\begin{align*}
    q = Q(\omega_u) = \frac{e^r - d}{u - d} \\
    1-q = Q(\omega_d) = \frac{u - e^r}{u - d}
\end{align*}
Since $\mathcal{M}^e(S)$ is a singleton, the market is complete.
\end{proposition}

\begin{proof}
(NA) $\iff \mathcal{M}^e(S) \neq \emptyset$. We seek $Q$ equivalent to $P$ such that $E_Q[\hat{S}_1^1] = \hat{S}_0^1$, where $\hat{S}$ is the discounted price.
\begin{equation}
E_Q\left[\frac{S_1^1}{S_1^0}\right] = \frac{S_0^1}{S_0^0} \iff q \frac{su}{e^r} + (1-q) \frac{sd}{e^r} = s
\end{equation}
Simplifying by $s$ and multiplying by $e^r$, we obtain:
\begin{equation}
q u + (1-q) d = e^r \iff q(u-d) = e^r - d
\end{equation}
This equation has a unique solution $q = \frac{e^r - d}{u - d}$. The equivalence conditions $q>0$ and $1-q>0$ are equivalent to $e^r-d>0$ and $u-e^r>0$, i.e., $d < e^r < u$.
\end{proof}

\subsection{Pricing and Hedging a Contingent Asset}
Consider a contingent asset $\hat{B}$ whose payoff at maturity is $\hat{B}(\omega_u) = \hat{B}_u$ and $\hat{B}(\omega_d) = \hat{B}_d$.
Since the market is complete, there exists a unique no-arbitrage price $A_0$ and a unique replication portfolio $(H_0^0, H_0^1)$.

\paragraph{No-Arbitrage Price:}
The price is given by the expectation of the discounted payoff under the martingale measure $Q$:
\begin{equation}
A_0 = E_Q\left[\frac{\hat{B}}{e^r}\right] = \frac{1}{e^r} \left( q \hat{B}_u + (1-q) \hat{B}_d \right) \quad \text{with} \quad q = \frac{e^r - d}{u - d}.
\end{equation}

\paragraph{Hedging Portfolio:}
The portfolio $(H_0^0, H_0^1)$ must satisfy $A_0 = H_0^0 S_0^0 + H_0^1 S_0^1 = H_0^0 + H_0^1 s$.
The portfolio value at time $T=1$ must replicate the payoff:
\begin{equation}
H_0^0 S_1^0 + H_0^1 S_1^1 = \hat{B}
\end{equation}
This leads to the system of equations:
\begin{equation}
\begin{cases}
    H_0^0 e^r + H_0^1 su = \hat{B}_u \\
    H_0^0 e^r + H_0^1 sd = \hat{B}_d
\end{cases}
\end{equation}
Subtracting the two lines gives $H_0^1 s(u-d) = \hat{B}_u - \hat{B}_d$. We deduce the quantity of risky asset to hold (called the "Delta"):
\begin{equation}
H_0^1 = \frac{\hat{B}_u - \hat{B}_d}{s(u-d)}
\end{equation}
And the quantity of risk-free asset:
\begin{equation}
H_0^0 = A_0 - H_0^1 s
\end{equation}
In the binomial model with $d < e^r < u$, I can explicitly price and hedge any contingent asset.

\section{The Trinomial Model}
We extend the previous model to a three-state universe $\Omega = \{\omega_u, \omega_m, \omega_d\}$. The final prices of the risky asset are $su, sm, sd$ with $u > m > d$.

\begin{proposition}
The no-arbitrage assumption (NA) is satisfied if and only if $d < e^r < u$.
\end{proposition}

\begin{proof}

(NA) $\iff \mathcal{M}^e(S) \neq \emptyset$. We seek a probability $Q=(q_1, q_2, q_3)$ with $q_i = Q(\omega_i) > 0$ such that:
\begin{align}
    q_1 + q_2 + q_3 &= 1 \quad (P_1) \text{ (probability condition)} \\
    q_1 u + q_2 m + q_3 d &= e^r \quad (P_2) \text{ (martingale condition)}
\end{align}

\begin{tcolorbox}[title=Geometric Visualization of Arbitrage, colback=yellow!10!white, colframe=yellow!50!black]
The set of probabilities is represented by a \textbf{simplex} (red triangle). For there to be no arbitrage (NA), the martingale plane $(P_2)$ must cross the interior of this triangle. This occurs if and only if $d < e^r < u$.
\end{tcolorbox}
\end{proof}

\paragraph{How to Read Figure 4.1:}
This figure illustrates the structure of the set of equivalent martingale measures $\mathcal{M}^e(S)$. Several key elements should be observed:
\begin{itemize}
    \item \textbf{The Simplex (Red):} It contains all probabilities $Q$ such that $\sum q_i = 1$. The interior represents \textit{equivalent} probabilities ($q_i > 0$).
    \item \textbf{The Martingale Set ($\mathcal{M}^e(S)$):} This is the line segment (blue or green) resulting from the intersection of the simplex with the martingale plane.
    \item \textbf{The Dashed Square (Gray):} On the horizontal plane ($q_1, q_3$), it delimits the domain $[0,1] \times [0,1]$. It represents the space where each individual probability is valid. The simplex only occupies the "lower" part ($q_1+q_3 \le 1$) because the sum of the three probabilities must be exactly 1.
    \item \textbf{Dependence on Parameters:} Depending on whether the "middle" return $m$ is greater or less than the risk-free rate $e^r$, the segment pivots to connect different axes of the simplex.
\end{itemize}

\begin{figure}[H]
\centering
\tdplotsetmaincoords{70}{120}
\begin{tikzpicture}[tdplot_main_coords, scale=4]
    \coordinate (O) at (0,0,0);
    
    % Axes setup:
    % q2 Up -> Z axis (0,0,1)
    % q1 Right -> Y axis (0,1,0)
    % q3 Down/Left -> X axis (1,0,0)
    
    \draw[-{Latex[length=2mm]}] (O) -- (1.5,0,0) node[anchor=north east]{$q_3$};
    \draw[-{Latex[length=2mm]}] (O) -- (0,1.5,0) node[anchor=north west]{$q_1$};
    \draw[-{Latex[length=2mm]}] (O) -- (0,0,1.5) node[anchor=south]{$q_2$};
    
    % Unit Square on floor (q1, q3)
    \draw[dashed, gray!80] (1,0,0) -- (1,1,0) -- (0,1,0);
    \node[gray, font=\tiny, anchor=south west] at (1,1,0) {$(1,1)$};

    % Vertices of Simplex on axes (adjusted to match professor's schema)
    \coordinate (Q3) at (1.25,0,0); % q3 axis intercept - outside unit square
    \coordinate (Q1) at (0,0.75,0); % q1 axis intercept - inside unit square
    \coordinate (Q2) at (0,0,1); % q2 axis intercept

    % Simplex (P1) - RED
    \draw[fill=red!5, opacity=0.6] (Q1) -- (Q2) -- (Q3) -- cycle;
    \draw[red, thick] (Q1) -- (Q2);
    \draw[red, thick] (Q2) -- (Q3);
    \draw[red, thick] (Q3) -- (Q1);
    
    \node[red, font=\bfseries] at (0.2, 0.2, 0.9) {Simplex $(P_1)$};

    % Blue Setup: Line on Q1-Q3 edge to Q2-Q3 edge
    % Point on q1-q3 edge (floor): parametric from (0,0.75,0) to (1.25,0,0)
    \coordinate (Blue_P1) at (0.9, 0.2, 0); 
    % Point on q2-q3 edge (x-z plane): parametric from (0,0,1) to (1.25,0,0)
    \coordinate (Blue_P2) at (0.85, 0, 0.3);
    
    \draw[ultra thick, blue] (Blue_P1) -- (Blue_P2) node[pos=0.1, left, xshift=-10pt, yshift=5pt, font=\scriptsize]{$\mathcal{M}^e(S)$};
    \fill[blue] (Blue_P1) circle (0.4pt);
    \fill[blue] (Blue_P2) circle (0.4pt);
    
    % Blue Normal n1
    \draw[->, blue, thick] (O) -- (1.5, 1.5, 1.5) node[above right, blue, inner sep=2pt] {$\vec{n}_1$};

    % Green Setup: Line on Q1-Q3 edge to Q1-Q2 edge
    % Point on q1-q3 edge (floor): closer to Q1
    \coordinate (Green_P1) at (0.4, 0.5, 0);
    % Point on q1-q2 edge (y-z plane): parametric from (0,0.75,0) to (0,0,1)
    \coordinate (Green_P2) at (0, 0.4, 0.45);

    \draw[ultra thick, green!60!black] (Green_P1) -- (Green_P2) node[pos=0.1, right, xshift=10pt, yshift=5pt, font=\scriptsize]{$\mathcal{M}^e(S)$};
    \fill[green!60!black] (Green_P1) circle (0.4pt);
    \fill[green!60!black] (Green_P2) circle (0.4pt);
    
    % Green Normal n2
    \draw[->, green!60!black, thick] (O) -- (0.2, 0.9, 0.4) node[right, green!60!black, inner sep=3pt] {$\vec{n}_2$};

    % Labels for points
    \node[below=8pt, blue] at (Blue_P1) {$\bar{P}_1$}; 
    \node[left=10pt, blue] at (Blue_P2) {$\bar{P}_2$};
    
    \node[below right=8pt, green!60!black] at (Green_P1) {$\bar{P}_1$}; 
    \node[right=10pt, green!60!black] at (Green_P2) {$\bar{P}_2'$};

\end{tikzpicture}
\caption{Geometry of incompleteness in the trinomial model. The intersection (thick segment) between the probability space and the martingale constraint shows the existence of infinitely many RNM measures. The blue configuration corresponds to $m \ge e^r$ and the green to $m < e^r$.}
\end{figure}

\begin{tcolorbox}[title=The Interest of the Geometric View, colback=red!5!white, colframe=red!50!black]
This view allows understanding that the price of an asset is no longer unique: it can vary along the segment $\mathcal{M}^e(S)$. The no-arbitrage price bounds are simply the asset values at the segment endpoints ($\bar{P}_1$ and $\bar{P}_2$).
\end{tcolorbox}

More generally, to obtain $\mathcal{M}^e(S)$, we seek the intersection between the interior of the simplex (open) and the plane $(P_2)$ with equation:
\begin{equation}
    q_1 u + q_2 m + q_3 d = e^r
\end{equation}
The following vector is orthogonal to plane $(P_2)$:
\begin{equation}
    \vec{V} = \begin{pmatrix} u/e^r \\ m/e^r \\ d/e^r \end{pmatrix}
\end{equation}
Under the no-arbitrage assumption (NA), we always have:
\begin{equation}
    \frac{u}{e^r} > 1 \quad \text{and} \quad 0 < \frac{d}{e^r} < 1
\end{equation}
The intersection then depends on the position of $\frac{m}{e^r}$ relative to 1. We distinguish two cases (illustrated in Figure 4.1):
To find the price bounds of a claim, it suffices to evaluate its expectation at the two extreme points of this segment, denoted $\bar{P}_1$ and $\bar{P}_2$.

The exact form of these extreme points depends on the position of the intermediate return $m$ relative to the risk-free rate $e^r$.
The set of martingale measures $\mathcal{M}^e(S)$ is the intersection of two geometric constraints: the probability simplex ($q_1+q_2+q_3=1, q_i>0$) and the martingale condition plane ($q_1u + q_2m + q_3d = e^r$). In the trinomial case with $d < e^r < u$, this intersection is a line segment. The no-arbitrage price bounds are attained at the endpoints of this segment. These extreme points correspond to martingale measures that "ignore" one of the three possible states, i.e., by assigning it zero probability ($q_i=0$).

Let us detail the calculation for the case where $m \ge e^r$.

\begin{itemize}
    \item \textbf{Calculation of $\bar{P}_1$:} This point is obtained by assuming that the intermediate state $m$ has zero probability, i.e., $q_2=0$. The system of equations reduces to:
    \begin{enumerate}
        \item $q_1 + q_3 = 1$
        \item $q_1u + q_3d = e^r$
    \end{enumerate}
    Substituting $q_3 = 1 - q_1$ into the second equation, we get $q_1u + (1 - q_1)d = e^r$, which gives $q_1(u - d) = e^r - d$. The solution is:
    \[
    q_1 = \frac{e^r - d}{u - d} \quad \text{and} \quad q_3 = 1 - q_1 = \frac{u - e^r}{u - d}.
    \]
    Thus $\bar{P}_1 = ( \frac{e^r-d}{u-d}, 0, \frac{u-e^r}{u-d} )$.

    \item \textbf{Calculation of $\bar{P}_2$:} This point is obtained by setting $q_1=0$. The system becomes:
    \begin{enumerate}
        \item $q_2 + q_3 = 1$
        \item $q_2m + q_3d = e^r$
    \end{enumerate}
    The resolution is analogous and gives:
    \[
    q_2 = \frac{e^r - d}{m - d} \quad \text{and} \quad q_3 = \frac{m - e^r}{m - d}.
    \]
    Thus $\bar{P}_2 = ( 0, \frac{e^r-d}{m-d}, \frac{m-e^r}{m-d} )$.
\end{itemize}

The calculation for the case $m < e^r$ follows the same logic, setting $q_2=0$ then $q_3=0$.

\subsection{Pricing of Contingent Assets}
For a contingent asset $\hat{B}$ with payoffs $(\hat{B}_u, \hat{B}_m, \hat{B}_d)$, the set of no-arbitrage prices is the open interval $]\underline{\Pi}(\hat{B}), \overline{\Pi}(\hat{B})[$, where:
\begin{itemize}
    \item $\underline{\Pi}(\hat{B}) = \inf_{Q \in \mathcal{M}^e(S)} E_Q[\hat{B}/e^r]$
    \item $\overline{\Pi}(\hat{B}) = \sup_{Q \in \mathcal{M}^e(S)} E_Q[\hat{B}/e^r]$
\end{itemize}
The map $Q \mapsto E_Q[\cdot]$ being linear and the set $\mathcal{M}^e(S)$ being convex (a segment), the supremum and infimum are attained at the extreme points of the segment. These extreme points, denoted $\bar{P}_1$ and $\bar{P}_2$, correspond to cases where one of the states has zero probability.

Two cases arise depending on the position of $m$ relative to $e^r$:
\begin{enumerate}
    \item If $m \ge e^r$: The extreme points are $\bar{P}_1$ (intersection with $q_2=0$) and $\bar{P}_2$ (intersection with $q_1=0$).
        \begin{itemize}
            \item $\bar{P}_1 = (\frac{e^r-d}{u-d}, 0, \frac{u-e^r}{u-d})$
            \item $\bar{P}_2 = (0, \frac{e^r-d}{m-d}, \frac{m-e^r}{m-d})$
        \end{itemize}
    \item If $m < e^r$: The extreme points are $\bar{P}_1$ (intersection with $q_2=0$) and $\bar{P}_2$ (intersection with $q_3=0$).
        \begin{itemize}
            \item $\bar{P}_1 = (\frac{e^r-d}{u-d}, 0, \frac{u-e^r}{u-d})$
            \item $\bar{P}_2 = (\frac{e^r-m}{u-m}, \frac{u-e^r}{u-m}, 0)$
        \end{itemize}
\end{enumerate}

\begin{example}[European Call Option]
Consider a call option with strike $K$ on the risky asset, so that the discounted payoff is $\hat{B} = (S_1^1 - K)^+/e^r$. Suppose $su > K > sm$ and $sd < K$. The payoffs are:
\[
\hat{B}(\omega_u) = \frac{su-K}{e^r}, \quad \hat{B}(\omega_m) = 0, \quad \hat{B}(\omega_d) = 0
\]
The price bounds are then:
\begin{itemize}
    \item $\overline{\Pi}(\hat{B}) = \max(E_{\bar{P}_1}[\hat{B}/e^r], E_{\bar{P}_2}[\hat{B}/e^r])$.
    Since $u$ is the only state where the payoff is non-zero, the sup is attained for the probability that maximizes $q_1$. This is $\bar{P}_1$ if $m < e^r$ or $\bar{P}_2$ if $m > e^r$ (depending on values).
    Suppose $m < e^r$. The max of $q_1$ is at $\bar{P}_2$.
    $\overline{\Pi}(\hat{B}) = E_{\bar{P}_2}[\hat{B}/e^r] = \frac{e^r-m}{u-m} \cdot \frac{su-K}{e^r}$.
    \item $\underline{\Pi}(\hat{B}) = \min(E_{\bar{P}_1}[\hat{B}/e^r], E_{\bar{P}_2}[\hat{B}/e^r])$.
    The min is at $\bar{P}_1$.
    $\underline{\Pi}(\hat{B}) = E_{\bar{P}_1}[\hat{B}/e^r] = \frac{e^r-d}{u-d} \cdot \frac{su-K}{e^r}$.
\end{itemize}
\end{example}

\begin{example}[Numerical Exercise - Trinomial Model]
Consider a one-period trinomial model with the following parameters: risk-free rate $r=0$ (hence $e^r = 1$), initial price $S_0^1 = 1$, return factors $u=1.2$, $m=1$, $d=0.8$. We consider a call option with strike $K = 0.9$, whose discounted payoff is $\hat{B} = (\hat{S}^1_1 - 0.9)^+$.

\textbf{Questions:}
\begin{enumerate}
    \item Verify the no-arbitrage condition.
    \item Calculate the payoffs in each state.
    \item Determine the extreme martingale measures.
    \item Deduce the no-arbitrage price interval.
    \item Construct a super-replication portfolio.
\end{enumerate}
\end{example}

\section*{Solution to Numerical Exercise 4.3}

This subsection provides a detailed solution to the above exercise.

\begin{enumerate}
    \item \textbf{Parameters:}
    \begin{itemize}
        \item Risk-free rate: $r=0$, hence $e^r = 1$.
        \item Initial price: $S_0^1 = 1$.
        \item Return factors: $u=1.2, m=1, d=0.8$.
        \item Contingent claim: $\hat{B} = (\hat{S}^1_1 - 0.9)^+$.
    \end{itemize}
    \item \textbf{Market Analysis:} We verify the no-arbitrage condition: $d < e^r < u$ becomes $0.8 < 1 < 1.2$. The condition is satisfied. The market is incomplete.
    \item \textbf{Payoff Calculation:}
    \begin{itemize}
        \item State $u$: $\hat{S}^1_1 = 1.2$, $\hat{B}_u = (1.2 - 0.9)^+ = 0.3$.
        \item State $m$: $\hat{S}^1_1 = 1.0$, $\hat{B}_m = (1.0 - 0.9)^+ = 0.1$.
        \item State $d$: $\hat{S}^1_1 = 0.8$, $\hat{B}_d = (0.8 - 0.9)^+ = 0$.
    \end{itemize}
    \item \textbf{Determination of Extreme Measures:} We are in the case $m = e^r = 1$, which is a limiting case of $m \ge e^r$.
    \begin{itemize}
        \item $\bar{P}_1$ (with $q_2=0$, hence $q_1+q_3=1$ and $1.2q_1 + 0.8q_3 = 1$):
        \[
        q_1 = \frac{1-0.8}{1.2-0.8} = \frac{0.2}{0.4} = 0.5.
        \]
        Thus $\bar{P}_1 = (0.5, 0, 0.5)$.
        \item $\bar{P}_2$ (with $q_1=0$, hence $q_2+q_3=1$ and $1.0q_2 + 0.8q_3 = 1$):
        \[
        q_2 = \frac{1-0.8}{1-0.8} = 1.
        \]
        Thus $\bar{P}_2 = (0, 1, 0)$.
    \end{itemize}
    \item \textbf{Price Bound Calculation:} The no-arbitrage price interval is $]\underline{\Pi}(\hat{B}), \overline{\Pi}(\hat{B})[$:
    \begin{itemize}
        \item $E_{\bar{P}_1}[\hat{B}] = 0.5 \times 0.3 + 0 \times 0.1 + 0.5 \times 0 = 0.15$.
        \item $E_{\bar{P}_2}[\hat{B}] = 0 \times 0.3 + 1 \times 0.1 + 0 \times 0 = 0.1$.
    \end{itemize}
    The price interval is therefore $]0.1, 0.15[$.
    \item \textbf{Super-Replication Portfolio Calculation:} The minimum cost to super-replicate is $\alpha = \overline{\Pi}(\hat{B}) = 0.15$. We need to find a portfolio $(H_0^0, H_0^1)$ such that its initial cost is 0.15 and its final value dominates the payoff $\hat{B}$ in all states.
    \begin{itemize}
        \item Initial cost: $\hat{V}_0 = H_0^0 \cdot 1 + H_0^1 \cdot 1 = 0.15$.
        \item The proposed portfolio is $H_0^1 = 0.75$ and $H_0^0 = -0.6$.
        \item Cost verification: $-0.6 + 0.75 = 0.15$. Correct.
        \item Super-replication verification at $T=1$:
        \begin{itemize}
            \item State $u$: $\hat{V}_1 = -0.6 \times 1 + 0.75 \times 1.2 = -0.6 + 0.9 = 0.3$. Now $\hat{B}_u = 0.3$. So $\hat{V}_1 \ge \hat{B}_u$. OK.
            \item State $m$: $\hat{V}_1 = -0.6 \times 1 + 0.75 \times 1.0 = 0.15$. Now $\hat{B}_m = 0.1$. So $\hat{V}_1 \ge \hat{B}_m$ ($0.15 \ge 0.1$). OK.
            \item State $d$: $\hat{V}_1 = -0.6 \times 1 + 0.75 \times 0.8 = -0.6 + 0.6 = 0$. Now $\hat{B}_d = 0$. So $\hat{V}_1 \ge \hat{B}_d$. OK.
        \end{itemize}
    \end{itemize}
    The portfolio successfully super-replicates the payoff for a cost of 0.15.
\end{enumerate}

% ========================================
% Chapter 5: Multi-Period Binomial Model and American Options
% ========================================
\chapter{Multi-Period Binomial Model and American Options}

\section{Multi-Period Binomial Model ($T=3$)}

Consider a binomial model with $T=3$ periods.  
The filtration is given by $(\mathcal{J}_t)_{t=0,\dots,3}$ with:
\[
\mathcal{J}_0=\{\emptyset,\Omega\},\quad
\mathcal{J}_1=\sigma(A_1,A_2),\quad
\mathcal{J}_2=\sigma(B_1,B_2,B_3,B_4),\quad
\mathcal{J}_3=\mathcal{P}(\Omega),
\]
where
\begin{itemize}
    \item $A_1 = \{\omega \in \Omega : \omega_1 = u\}$,
    \item $A_2 = \{\omega \in \Omega : \omega_1 = d\}$,
    \item $B_1 = \{\omega : \omega_1 = u,\ \omega_2 = u\}$,
    \item $B_2 = \{\omega : \omega_1 = u,\ \omega_2 = d\}$,
    \item $B_3 = \{\omega : \omega_1 = d,\ \omega_2 = u\}$,
    \item $B_4 = \{\omega : \omega_1 = d,\ \omega_2 = d\}$.
\end{itemize}

We consider a risk-free asset (bond) with price $S_t^0 = e^{rt}$  
and a risky asset (stock) with price $S_t^1$.

The evolution of the stock price can be represented by a binomial tree:

\begin{figure}[H]
\centering
\begin{tikzpicture}[
    grow=right,
    scale=0.85, transform shape,
    level distance=45mm,
    level 1/.style={sibling distance=35mm},
    level 2/.style={sibling distance=18mm},
    level 3/.style={sibling distance=9mm},
    every node/.style={align=center, font=\footnotesize},
    edge from parent/.style={draw, -latex, thin},
    edge label/.style={font=\scriptsize, inner sep=1pt, sloped, above}
]

\node {$S_0$}
    child {
        node {$S_1(d)=Sd$}
        child {
            node {$S_2(dd)=Sd^2$}
            child {
                node {$S_3(ddd)=Sd^3$}
                edge from parent node[below] {$1-p$}
            }
            child {
                node {$S_3(ddu)=Sd^2u$}
                edge from parent node[above] {$p$}
            }
            edge from parent node[below] {$1-p$}
        }
        child {
            node {$S_2(du)=Sdu$}
            child {
                node {$S_3(dud)=Sdud$}
                edge from parent node[below] {$1-p$}
            }
            child {
                node {$S_3(duu)=Sdu^2$}
                edge from parent node[above] {$p$}
            }
            edge from parent node[above] {$p$}
        }
        edge from parent node[below] {$1-p$}
    }
    child {
        node {$S_1(u)=Su$}
        child {
            node {$S_2(ud)=Sud$}
            child {
                node {$S_3(udd)=Sud^2$}
                edge from parent node[below] {$1-p$}
            }
            child {
                node {$S_3(udu)=Sudu$}
                edge from parent node[above] {$p$}
            }
            edge from parent node[below] {$1-p$}
        }
        child {
            node {$S_2(uu)=Su^2$}
            child {
                node {$S_3(uud)=Su^2d$}
                edge from parent node[below] {$1-p$}
            }
            child {
                node {$S_3(uuu)=Su^3$}
                edge from parent node[above] {$p$}
            }
            edge from parent node[above] {$p$}
        }
        edge from parent node[above] {$p$}
    };
\end{tikzpicture}

\caption{Evolution of stock price in a 3-period binomial model.}
\end{figure}

\section{Absence of Arbitrage and Risk-Neutral Measure}

The no-arbitrage condition (NA) in the $T=3$ model is equivalent to its validity in each one-period sub-model.  
In the binomial case, it reads:
\begin{equation}
d < e^r < u.
\end{equation}

The existence of an equivalent martingale measure (EMM) $Q$ is equivalent to (NA).  
Under $Q$, the discounted stock price is a martingale:
\begin{equation}
\E_Q[\hat{S}_{t+1} \mid \mathcal{J}_t] = \hat{S}_t,
\end{equation}
where $\hat{S}_t = S_t / S_t^0 = e^{-rt} S_t$.

For any node $(\omega_1, \dots, \omega_t)$, we have:
\[
e^{-rt} S_t(\omega_1, \dots, \omega_t)
= \E_Q[e^{-r(t+1)} S_{t+1} \mid \mathcal{J}_t],
\]
that is:
\[
S_t = e^{-r} \left[ S_{t+1}(u)\, Q_{(\omega_1,\dots,\omega_t)}(u)
+ S_{t+1}(d)\, Q_{(\omega_1,\dots,\omega_t)}(d) \right].
\]

Let us solve for a particular node, for example $t=2$ after two up moves:
\[
S_2(uu) = e^{-r} \big[S_3(uuu)\, q_{uu} + S_3(uud)\,(1-q_{uu})\big],
\]
that is
\[
S u^2 = e^{-r} \big[S u^3 q_{uu} + S u^2 d (1-q_{uu})\big].
\]
We then obtain
\[
1 = e^{-r}\,[u q_{uu} + d(1-q_{uu})]
\quad\Rightarrow\quad
e^r = d + (u-d) q_{uu},
\]
hence the risk-neutral probability:
\[
q_{uu} = q = \frac{e^r - d}{u - d}, \qquad
1-q_{uu} = \frac{u - e^r}{u - d}.
\]
In this model, returns are path-independent, so the probability $q$ is the same at each node.

\section{Valuation and Hedging of a Contingent Claim}

To value a contingent claim with payoff $B$ at maturity $T=3$, we use the risk-neutral valuation formula:
\begin{equation}
A_t = \E_Q[B \mid \mathcal{J}_t],
\end{equation}
where $A_t$ denotes the claim price at date $t$.

To hedge the claim, we construct a self-financing portfolio $(H_t^0,H_t^1)$ such that its value coincides with the claim value at each date.  
The portfolio composition is given by:
\begin{align}
H_t^1 &= \frac{A_{t+1}(u) - A_{t+1}(d)}{S_{t+1}(u) - S_{t+1}(d)}
\quad &\text{(number of shares held)}\\
H_t^0 &= e^{-r}\,\big(A_{t+1}(u) - H_t^1 S_{t+1}(u)\big)
\quad &\text{(amount invested in risk-free asset).}
\end{align}
Since the market is complete, the measure $Q$ is unique and any claim can be perfectly replicated.

\section*{Numerical Example: Binomial Model}

Parameters: $T=3$, $u=1.1$, $d=0.9$, $r=\log(1+0.05)$, $S_0=20$

\subsection*{Stock Price Tree}

\begin{center}
\begin{tikzpicture}[scale=0.85, transform shape]
    % Node styles
    \tikzstyle{n} = [align=center, font=\small]
    
    % Nodes
    % t=0
    \node[n] (n0) at (0,0) {20};
    
    % t=1
    \node[n] (n1u) at (3, 1.5) {22 \\ \textcolor{red}{(20.95)}};
    \node[n] (n1d) at (3, -1.5) {18 \\ \textcolor{red}{(17.14)}};
    
    % t=2
    \node[n] (n2uu) at (6, 3) {24.2 \\ \textcolor{red}{(21.95)}};
    \node[n] (n2ud) at (6, 0) {19.8 \\ \textcolor{red}{(17.96)}};
    \node[n] (n2dd) at (6, -3) {16.2 \\ \textcolor{red}{(14.69)}};
    
    % t=3
    \node[n] (n3uuu) at (9, 4.5) {26.62 \\ \textcolor{red}{(23)}};
    \node[n] (n3uud) at (9, 1.5) {21.78 \\ \textcolor{red}{(18.81)}};
    \node[n] (n3udd) at (9, -1.5) {17.82 \\ \textcolor{red}{(15.39)}};
    \node[n] (n3ddd) at (9, -4.5) {14.58 \\ \textcolor{red}{(12.59)}};
    
    % Edges
    \draw[->] (n0) -- (n1u);
    \draw[->] (n0) -- (n1d);
    
    \draw[->] (n1u) -- (n2uu);
    \draw[->] (n1u) -- (n2ud);
    \draw[->] (n1d) -- (n2ud);
    \draw[->] (n1d) -- (n2dd);
    
    \draw[->] (n2uu) -- (n3uuu);
    \draw[->] (n2uu) -- (n3uud);
    \draw[->] (n2ud) -- (n3uud);
    \draw[->] (n2ud) -- (n3udd);
    \draw[->] (n2dd) -- (n3udd);
    \draw[->] (n2dd) -- (n3ddd);

\end{tikzpicture}
\end{center}

\vspace{1cm}

\noindent Notation: $\hat{S}_t$ and $\left(\hat{S}_t^1\right)$

\vspace{1cm}

\subsection*{Price of a European Put with $K=20$ and $T=3$}

$$\hat{B} = \left(-\hat{S}_T + K\right)_+ \quad \text{with} \quad K \cdot e^{-3r} \approx 17.28 \text{ (discounted strike value)}$$

\newpage

\subsection*{Risk-Neutral Probabilities}

Here $Q_{(w_1,w_2)}(w_u) = Q_{w_1}(w_u) = Q(w_u) = \frac{3}{4}$

$Q_{(w_1,w_2)}(w_d) = Q_{w_1}(w_d) = Q(w_d) = \frac{1}{4}$

\vspace{0.5cm}

$$B = \left(\frac{K}{e^{3r}} - S_3^1\right)_+ = A_3 = \begin{cases}
0 \\
0 \\
1.89 \\
4.69
\end{cases}$$

where $e^{3r} = 17.28$

\subsection*{Tree for $A_t$ (Option Price)}

The following binomial tree represents the European put option price at each node. Nodes are numbered from \circled{1} to \circled{6} and the discounted terminal payoff values ($A_3$) are indicated on the right.

\begin{center}
\begin{center}
\begin{tikzpicture}[scale=0.85, transform shape]
    % Node styles
    \tikzstyle{n} = [circle, draw, minimum size=0.9cm, font=\small, fill=white]
    \tikzstyle{l} = [font=\small, align=left]

    % Coordinates for Recombining Tree
    % t=0
    \node[n] (n1) at (0,0) {\circled{1}};
    
    % t=1
    \node[n] (n3) at (3, 1.5) {\circled{3}};
    \node[n] (n2) at (3, -1.5) {\circled{2}};
    
    % t=2
    \node[n] (n6) at (6, 3) {\circled{6}};
    \node[n] (n5) at (6, 0) {\circled{5}};
    \node[n] (n4) at (6, -3) {\circled{4}};
    
    % t=3 (Leaves/Payoffs)
    \node[right] at (8, 4.5) {0};
    \node[right] at (8, 1.5) {0};
    \node[right] at (8, -1.5) {1.89};
    \node[right] at (8, -4.5) {4.69};
    
    % Edges
    % t=0 -> t=1
    \draw[->, thick] (n1) -- (n3);
    \draw[->, thick] (n1) -- (n2);
    
    % t=1 -> t=2
    \draw[->, thick] (n3) -- (n6);
    \draw[->, thick] (n3) -- (n5);
    \draw[->, thick] (n2) -- (n5);
    \draw[->, thick] (n2) -- (n4);
    
    % t=2 -> t=3 (Payoff links)
    \draw[->] (n6) -- (7.5, 4.5);
    \draw[->] (n6) -- (7.5, 1.5);
    \draw[->] (n5) -- (7.5, 1.5);
    \draw[->] (n5) -- (7.5, -1.5);
    \draw[->] (n4) -- (7.5, -1.5);
    \draw[->] (n4) -- (7.5, -4.5);

\end{tikzpicture}
\end{center}
\end{center}

\begin{center}
\underline{Discounted Payoff}
\end{center}

\textbf{Backward Calculation of $A_t$:} We work backwards through the tree calculating the discounted expectation under $Q$ (with $q = 3/4$):

\begin{align*}
\circled{6} &= q \times 0 + (1-q) \times 0 = 0 \\
\circled{5} &= q \times 0 + (1-q) \times 1.89 \approx 0.47 \\
\circled{4} &= q \times 1.89 + (1-q) \times 4.69 \approx 2.58 \\[1em]
\circled{3} &= q \times \circled{6} + (1-q) \times \circled{5} \approx 0.12 \\
\circled{2} &= q \times \circled{5} + (1-q) \times \circled{4} \approx 1.00 \\[1em]
\circled{1} &= q \times \circled{3} + (1-q) \times \circled{2} \approx 0.34
\end{align*}

\newpage

\subsection*{Tree for $\textcolor{teal}{H_t^0}$ and $\textcolor{red}{H_t^1}$}

The following tree shows the hedging strategy values at each node. The \textcolor{teal}{teal} nodes represent $H^0$ (risk-free asset position) and the \textcolor{red}{red} nodes represent $H^1$ (Delta, risky asset position).

\begin{center}
\begin{center}
\begin{center}
\begin{tikzpicture}[scale=0.85, transform shape]
    % Node styles
    % Double circle to represent both H^0 (teal) and H^1 (red)
    \tikzstyle{n} = [circle, draw=teal, double=red, double distance=1.5pt, minimum size=0.9cm, font=\small, fill=white, thick]
    
    % Tree Layout (Non-recombining for Strategy as per notes)
    % t=0
    \node[n] (n6) at (0,0) {\circled{6}};
    
    % t=1
    \node[n] (n4) at (3, 1.5) {\circled{4}};
    \node[n] (n5) at (3, -1.5) {\circled{5}};
    
    % t=2
    \node[n] (n1) at (6, 2.5) {\circled{1}}; % UU
    \node[n] (n2) at (6, 0) {\circled{2}}; % UD/DU (Recombining)
    \node[n] (n3) at (6, -2.5) {\circled{3}}; % DD
    
    % Edges
    \draw[->, thick, gray] (n6) -- (n4);
    \draw[->, thick, gray] (n6) -- (n5);
    
    \draw[->, thick, gray] (n4) -- (n1);
    \draw[->, thick, gray] (n4) -- (n2);
    
    \draw[->, thick, gray] (n5) -- (n2);
    \draw[->, thick, gray] (n5) -- (n3);

\end{tikzpicture}
\end{center}
\end{center}
\end{center}

\subsection*{For Example}

\textbf{Calculation of $H_2^1(w_u, w_u)$:}
\[
\circled{3} = H_2^1(w_u, w_u) = \frac{A_3(w_u, w_u, w_u) - A_3(w_u, w_u, w_d)}{S_3^1(w_u, w_u, w_u) - S_3^1(w_u, w_u, w_d)} = -1
\]

\textbf{Calculation of $H_2^0(w_u, w_u)$:} (financing asset)
\[
\circled{3} = H_2^0(w_u, w_u) = A_2(w_u, w_u) - H_2^1(w_u, w_u) \times S_2^1(w_u, w_u) = 17.27
\]

\subsection*{Similarly}

\begin{minipage}{0.45\textwidth}
\textbf{Values of $\textcolor{teal}{H^0}$:}
\begin{align*}
\textcolor{teal}{\circled{1}} &= 0 \\
\textcolor{teal}{\circled{2}} &= 10.35 \\
\textcolor{teal}{\circled{4}} &= 2.634 \\
\textcolor{teal}{\circled{5}} &= 11.95 \\
\textcolor{teal}{\circled{6}} &= 4.44
\end{align*}
\end{minipage}
\hfill
\begin{minipage}{0.45\textwidth}
\textbf{Values of $\textcolor{red}{H^1}$:}
\begin{align*}
\textcolor{red}{\circled{1}} &= 0 \\
\textcolor{red}{\circled{2}} &= -0.55 \\
\textcolor{red}{\circled{3}} &= -1 \\
\textcolor{red}{\circled{4}} &= -0.12 \\
\textcolor{red}{\circled{5}} &= -0.64 \\
\textcolor{red}{\circled{6}} &= -0.23
\end{align*}
\end{minipage}

\section{American Options}

\subsection{Definition}

An American option is a non-negative stochastic process $(\hat{X}_t)_{t=0,\dots,T}$ adapted to the filtration $(\mathcal{J}_t)$, which the holder can exercise at any time $t\in\{0,\dots,T\}$ to receive the payment $\hat{X}_t$.

\begin{itemize}
    \item American call option: $\hat{X}_t = (S_t - K)_+$,
    \item American put option: $\hat{X}_t = (K - S_t)_+$.
\end{itemize}

In a complete market, one cannot perfectly replicate an American option by a static strategy, because the exercise date is random.

\subsection{Valuation by Backward Induction}

Let $(\hat{U}_t)_{t=0,\dots,T}$ be the price process of an American option.  
At maturity: $\hat{U}_T = \hat{X}_T$.  
At date $T-1$, the holder chooses between:
\begin{enumerate}
    \item exercising and receiving $\hat{X}_{T-1}$;
    \item not exercising and keeping an option worth
    $\E_Q[e^{-r}\hat{X}_T \mid \mathcal{J}_{T-1}]$.
\end{enumerate}
Thus:
\[
\hat{U}_{T-1} = \max\!\left(\hat{X}_{T-1},\, \E_Q[e^{-r}\hat{U}_T\mid\mathcal{J}_{T-1}]\right),
\]
and, by induction, for all $t<T$:
\[
\hat{U}_t = \max\!\left(\hat{X}_t,\, \E_Q[e^{-r}\hat{U}_{t+1}\mid\mathcal{J}_t]\right).
\]

Setting $U_t = e^{-rt}\hat{U}_t$ and $X_t = e^{-rt}\hat{X}_t$ (discounted values), we obtain:
\begin{equation}
U_t = \max\!\left(X_t,\, \E_Q[U_{t+1}\mid\mathcal{J}_t]\right).
\end{equation}

\begin{proposition}[Snell Envelope]
The discounted process $(U_t)$ is the smallest supermartingale under $Q$ dominating the process $(X_t)$.
\end{proposition}

\begin{proof}
The proof proceeds in two steps:
\begin{enumerate}
    \item \textbf{$U$ dominates $X$ and is a supermartingale:}
    By definition, $U_t = \max(X_t, \E_Q[U_{t+1}|\mathcal{J}_t])$.
    Therefore $U_t \ge X_t$ (dominance).
    Moreover, $U_t \ge \E_Q[U_{t+1}|\mathcal{J}_t]$ (supermartingale).
    
    \item \textbf{$U$ is the smallest:}
    Let $(Y_t)$ be another supermartingale dominating $(X_t)$.
    At $t=T$, $U_T = X_T \le Y_T$.
    Suppose by induction that $U_{t+1} \le Y_{t+1}$.
    Then, by the supermartingale property of $Y$:
    \[
    \E_Q[U_{t+1}|\mathcal{J}_t] \le \E_Q[Y_{t+1}|\mathcal{J}_t] \le Y_t.
    \]
    Moreover, by the dominance hypothesis, $X_t \le Y_t$.
    Consequently:
    \[
    U_t = \max(X_t, \E_Q[U_{t+1}|\mathcal{J}_t]) \le \max(Y_t, Y_t) = Y_t.
    \]
    Thus, $U_t \le Y_t$ for all $t$, which proves that $U$ is the smallest dominating supermartingale.
\end{enumerate}
\end{proof}

\begin{proposition}[Non-Optimality of Early Exercise for a Call]
If the risk-free rate is non-negative ($S_t^0$ increasing), it is never optimal to exercise an American call option on a non-dividend-paying stock before maturity.
\end{proposition}

\begin{proof}[Detailed Proof]
We seek to prove that it is never optimal to exercise an American call option before its maturity $T$, provided the interest rate $r$ is positive. We reason by contradiction.

\begin{enumerate}
    \item \textbf{Contradiction hypothesis:} Suppose there exists a time $t < T$ where exercise is optimal. This means that the value obtained by exercising immediately, $S_t - K$ (intrinsic value), is strictly greater than the expected value if one keeps the option, $\E_Q[e^{-r} \hat{U}_{t+1} \mid \mathcal{J}_t]$ (continuation value).
    \[
    S_t - K > \E_Q[e^{-r} \hat{U}_{t+1} \mid \mathcal{J}_t]
    \]
    
    \item \textbf{Lower bound on future value:} The future option value, $\hat{U}_{t+1}$, is always at least equal to its intrinsic exercise value $(S_{t+1} - K)_+$.
    \[
    \hat{U}_{t+1} \ge (S_{t+1} - K)_+ \ge S_{t+1} - K
    \]
    
    \item \textbf{Taking expectation:} Take the discounted conditional expectation of this inequality:
    \[
    \E_Q[e^{-r} \hat{U}_{t+1} \mid \mathcal{J}_t] \ge \E_Q[e^{-r} (S_{t+1} - K) \mid \mathcal{J}_t]
    \]
    
    \item \textbf{Martingale Property:} The right-hand term expands by linearity:
    \[
    \E_Q[e^{-r} S_{t+1} \mid \mathcal{J}_t] - e^{-r} K
    \]
    Under measure $Q$, the discounted price $e^{-rt} S_t$ is a martingale. Therefore $\E_Q[e^{-r(t+1)}S_{t+1} \mid \mathcal{J}_t] = e^{-rt}S_t$. Multiplying by $e^{rt}$, we have $\E_Q[e^{-r} S_{t+1} \mid \mathcal{J}_t] = S_t$.
    The inequality thus becomes:
    \[
    \E_Q[e^{-r} \hat{U}_{t+1} \mid \mathcal{J}_t] \ge S_t - e^{-r} K
    \]
    
    \item \textbf{Contradiction:} Combine the initial hypothesis (1) and result (4):
    \[
    S_t - K > \E_Q[e^{-r} \hat{U}_{t+1} \mid \mathcal{J}_t] \ge S_t - e^{-r} K
    \]
    Simplifying by $S_t$, we obtain:
    \[
    -K > -e^{-r} K
    \]
    Multiplying by $-1$ (which reverses the inequality):
    \[
    K < e^{-r} K \iff K(1 - e^{-r}) < 0
    \]
    Now, if $r > 0$, then $e^{-r} < 1$, so $(1 - e^{-r}) > 0$. Since $K > 0$ as well, the product $K(1 - e^{-r})$ is \textbf{strictly positive}, which contradicts the inequality $K(1 - e^{-r}) < 0$.
    
    This contradiction shows that the initial hypothesis is false: it is never optimal to exercise an American call prematurely.
\end{enumerate}
\end{proof}

\subsection{Numerical Example: American Put Option}
Let us revisit the same binomial model as before: $S_0=20, u=1.1, d=0.9, r=\ln(1.05)$, but with an American put option with strike $K=22$.

\begin{tcolorbox}[title=Decision Tree: Exercise or Wait?, colback=blue!5!white, colframe=blue!50!black]
Calculations are done backward. At each node, we compare the intrinsic value $K-S_t$ with the continuation value.
\begin{itemize}
    \item $\mathbf{t=3}$: Payoff $P_3 = (22 - S_3)_+$.
    \begin{itemize}
        \item $S_3(ddd) = 14.58 \implies P_3 = 7.42$
        \item $S_3(udd) = 17.82 \implies P_3 = 4.18$
        \item $S_3(uud) = 21.78 \implies P_3 = 0.22$
        \item $S_3(uuu) = 26.62 \implies P_3 = 0$
    \end{itemize}
    \item $\mathbf{t=2}$: $P_2 = \max(22-S_2, e^{-r}[q P_3(u) + (1-q) P_3(d)])$.
    \begin{itemize}
        \item Node $dd$ ($S_2=16.2$): Intrinsic $= 5.8$. Continuation $\approx 4.75$. $\to$ \textbf{Exercise} (value 5.8).
        \item Node $ud$ ($S_2=19.8$): Intrinsic $= 2.2$. Continuation $\approx 1.15$. $\to$ \textbf{Exercise} (value 2.2).
        \item Node $uu$ ($S_2=24.2$): Intrinsic $= 0$. Continuation $\approx 0$. $\to$ No exercise.
    \end{itemize}
    \item $\mathbf{t=1}$:
    \begin{itemize}
        \item Node $d$ ($S_1=18$): Intrinsic $= 4$. Continuation $\approx 2.95$. $\to$ \textbf{Exercise} (value 4).
        \item Node $u$ ($S_1=22$): Intrinsic $= 0$. Continuation $\approx 0.52$. $\to$ No exercise.
    \end{itemize}
    \item $\mathbf{t=0}$ ($S_0=20$): Intrinsic $= 2$. Continuation $\approx 1.32$. $\to$ \textbf{Immediate Exercise} (value 2).
\end{itemize}
\end{tcolorbox}
Conclusion: In this specific example, the American put option has a value of 2 (immediate exercise), significantly higher than the corresponding European option (which would be worth approximately 1.32). This illustrates the early exercise premium.

\section{Stopping Times and American Options}

\begin{definition}[Stopping Time]
A random variable $\tau : \Omega \to \{0,\dots,T\}$ is a \emph{stopping time} if,
for all $t$, the event $\{\tau = t\}$ belongs to $\mathcal{J}_t$.
\end{definition}

\paragraph*{Intuition: The Best Time to Act}
The problem of American option valuation is a canonical example of an "optimal stopping problem". The option holder faces a series of decisions: at each time $t$, should they exercise the option and collect the gain $X_t$, or should they wait hoping for a higher future gain? A "stopping time" $\tau$ is a decision rule that depends only on the past and present. The holder seeks the optimal stopping time $\tau^*$ that maximizes the expected discounted gain, i.e., $\max_\tau \E_Q[X_\tau]$.

The backward induction formula we established ($U_t = \max(X_t, \E_Q[U_{t+1} \mid \mathcal{J}_t])$) is none other than Bellman's dynamic programming algorithm that solves this optimal stopping problem. At each step, it compares the value of immediate action ("exercise now") with the expected value of the optimal future strategy ("continue and make the best decision later"). The process $U_t$ thus represents the maximum attainable value from time $t$ onward.

Exercising an American option amounts to choosing a stopping time $\tau$.  
The gain obtained is $X_\tau$.  
\begin{proof}
We must show that the process defined by $U_t = \sup_{\tau \in \mathcal{T}_t} \E_Q[X_\tau | \mathcal{J}_t]$ indeed satisfies Bellman's recurrence equation (Proposition 1).
Let $V_t$ be this supremum.
At $t=T$, $V_T = X_T = U_T$.
For $t < T$, separate the set of stopping times $\mathcal{T}_t$ into two cases: $\tau = t$ (immediate exercise) and $\tau > t$ (continuation).
If $\tau > t$, then $\tau$ is also a stopping time in $\mathcal{T}_{t+1}$.
\[
V_t = \max \left( X_t, \sup_{\tau \in \mathcal{T}_{t+1}} \E_Q[X_\tau | \mathcal{J}_t] \right)
\]
Using the tower property (law of iterated expectations):
\[
\E_Q[X_\tau | \mathcal{J}_t] = \E_Q [ \E_Q[X_\tau | \mathcal{J}_{t+1}] | \mathcal{J}_t ]
\]
Therefore:
\[
\sup_{\tau \in \mathcal{T}_{t+1}} \E_Q[X_\tau | \mathcal{J}_t] = \E_Q \left[ \sup_{\tau \in \mathcal{T}_{t+1}} \E_Q[X_\tau | \mathcal{J}_{t+1}] \mid \mathcal{J}_t \right] = \E_Q[V_{t+1} | \mathcal{J}_t]
\]
We thus recover:
\[
V_t = \max(X_t, \E_Q[V_{t+1} | \mathcal{J}_t])
\]
This is exactly the recurrence relation that defines $(U_t)$. By uniqueness, $U_t$ is indeed the supremum of expected gains over all possible stopping times.
\end{proof}


\appendix

\end{document}
%===========================
